\documentclass[11pt]{article}

\usepackage[a4paper,margin=1in]{geometry}
\usepackage{amsmath,amssymb,amsthm,mathtools}
\usepackage{bm}
\usepackage{booktabs}
\usepackage{enumitem}
\usepackage{microtype}
\usepackage[authoryear,round]{natbib}
\setcitestyle{authoryear,open={(},close={)},aysep={ },citesep={; }}
\usepackage[colorlinks=true,allcolors=blue]{hyperref}
\usepackage[nameinlink,capitalise,noabbrev]{cleveref}
\crefname{assumption}{Assumption}{Assumptions}
\Crefname{assumption}{Assumption}{Assumptions}
\crefname{lemma}{Lemma}{Lemmas}
\Crefname{lemma}{Lemma}{Lemmas}
\crefname{proposition}{Proposition}{Propositions}
\Crefname{proposition}{Proposition}{Propositions}
\crefname{remark}{Remark}{Remarks}
\Crefname{remark}{Remark}{Remarks}
\crefname{definition}{Definition}{Definitions}
\Crefname{definition}{Definition}{Definitions}

\newtheorem{theorem}{Theorem}[section]
\newtheorem{proposition}[theorem]{Proposition}
\newtheorem{lemma}[theorem]{Lemma}
\newtheorem{corollary}[theorem]{Corollary}
\newtheorem{assumption}[theorem]{Assumption}
\newtheorem{remark}[theorem]{Remark}
\theoremstyle{definition}
\newtheorem{definition}[theorem]{Definition}

\newcommand{\R}{\mathbb{R}}
\newcommand{\E}{\mathbb{E}}
\newcommand{\Pp}{\mathbb{P}}
\newcommand{\Law}{\mathcal{L}}
\newcommand{\N}{\mathcal{N}}
\newcommand{\Id}{I_d}
\newcommand{\pdata}{p_{\mathrm{data}}}
\newcommand{\Wtwo}{W_2}
\newcommand{\norm}[1]{\left\lVert #1\right\rVert_2}
\newcommand{\ip}[2]{\left\langle #1,#2\right\rangle}
\newcommand{\Tr}{\operatorname{Tr}}
\newcommand{\Cov}{\operatorname{Cov}}
\newcommand{\supp}{\operatorname{supp}}
\newcommand{\KL}{\operatorname{KL}}
\newcommand{\TV}{\operatorname{TV}}
\newcommand{\scoreerr}{\varepsilon_{\mathrm{score}}}
\newcommand{\dd}{\,\mathrm{d}}
\newcommand{\lesssimc}{\lesssim}

\title{Intrinsic-Dimension-Dependent Wasserstein Bounds\\
for Denoising Diffusion Probabilistic Models}
\author{Kazuma Hiro}
\date{July 2026}

\begin{document}
\maketitle

\begin{abstract}
We study the convergence of the denoising diffusion probabilistic model (DDPM) when the target distribution has low intrinsic complexity. The data distribution may be singular with respect to Lebesgue measure; its low-dimensional structure is quantified by a metric-entropy condition on its support. For the 2-Wasserstein analysis, we represent the original DDPM update as the grid-point values of a semilinear stochastic differential equation and compare this interpolation with the exact reverse-time Ornstein--Uhlenbeck diffusion by a synchronous coupling. Tweedie's formula rewrites the nonlinear drift in terms of the posterior mean of the clean data. Along the reverse process, this posterior mean is a square-integrable martingale, which yields an orthogonal decomposition between the propagated sampling error and the local time-discretization error. We then refine the discretization analysis by expressing the posterior-covariance decrement as a single Lebesgue--Stieltjes integral with respect to the signal-to-noise ratio. The I--MMSE identity turns this integral into a mutual-information budget, and metric entropy bounds that budget by $k\log k$, where $k$ is the intrinsic dimension. Besides the resulting Wasserstein bound, we obtain an improved KL-divergence guarantee and reduce the sufficient iteration complexity from order $k\log^3(k/\varepsilon)/\varepsilon^2$ to $k\log^2(k/\varepsilon)/\varepsilon^2$.
\end{abstract}

\paragraph{Keywords.}
Diffusion models; DDPM; Wasserstein distance; KL divergence; intrinsic dimension; metric entropy; I--MMSE; reverse-time diffusion.

\section{Introduction}

Diffusion models have become a leading approach to high-dimensional
generative modeling \citep{ho2020ddpm,song2021score}. Given
independent observations from an unknown probability distribution
$\pdata$ on $\R^d$, a diffusion model gradually perturbs the data by
Gaussian noise and then learns to reverse this noising mechanism. If
$\hat p$ denotes the distribution produced by the resulting sampler,
a central theoretical problem is to quantify the discrepancy between
$\pdata$ and $\hat p$ under a probability metric or divergence.

\paragraph{Convergence analysis in KL divergence and total variation.}
A large part of the existing theory measures the sampling error by Kullback--Leibler divergence. 
This choice is analytically convenient because Girsanov's theorem converts 
an error between reverse-time drifts into a path-space relative-entropy estimate. 
Under weak assumptions on the data distribution, 
dimension-dependent convergence guarantees have been obtained by, among others, 
\citet{chen2023sampling}, \citet{chen2023improved}, \citet{benton2024nearly}, and \citet{koike2026follmer}. 
In particular, the sharpest general results exhibit a leading ambient-dimensional dependence of order $d$, up to the accuracy, score-estimation, and time-discretization factors appearing in the corresponding theorem; this order is essentially unavoidable without additional structure \citep{benton2024nearly}. Thus, from the viewpoint of KL divergence, the full-dimensional convergence theory of diffusion models is now comparatively well developed.

These KL guarantees also imply convergence in total variation through
Pinsker's inequality,
\[
    \TV(P,Q)
    \le
    \sqrt{\frac{1}{2}\KL(P\Vert Q)}.
\]
This implication is useful but takes a square root of the KL bound:
for example, a KL rate of order $d/N$ yields only a TV rate of order
$\sqrt{d/N}$, where $N$ is the number of sampling steps.
\citet[Theorem~1]{li2025od} instead analyze the DDPM recursion
directly in total variation. Under a finite first-moment condition and
$L_2$-accurate score estimates, their direct analysis gives, with
exact scores, a TV rate of order $\widetilde O(d/N)$, where
$\widetilde O$ suppresses logarithmic factors. Score-estimation error
enters their bound additively, up to logarithmic factors. Thus, their
result improves on the rate
$O(\sqrt{d/N})$ obtained by applying Pinsker's inequality to an
order-$d/N$ KL bound.

\paragraph{Convergence analysis in 2-Wasserstein distance.}
Convergence in the 2-Wasserstein distance has also received increasing attention \citep{gao2025wasserstein,bruno2025wasserstein,arsenyan2025assessing,gentiloni2025beyond,beyler2025convergence,yu2025advancing}. Existing results obtain the natural ambient-dimensional scale $\Wtwo(\pdata,\hat p)=O(\sqrt d)$ in several regimes. Compared with the KL theory, however, these results generally require stronger assumptions, such as log-concavity, semiconvexity, one-sided curvature conditions, or regularity of the true or learned score. The reason is that the path-space change-of-measure argument available for KL divergence does not directly control a Wasserstein distance; one must instead construct and stabilize a coupling between the exact and approximate sampling trajectories.

Wasserstein error is nevertheless important for both theoretical and practical reasons. It is defined for singular target distributions as soon as the relevant moments are finite, and it incorporates the geometry of the sample space through transportation cost. Moreover, the Fr\'echet Inception Distance used in image-generation evaluation is the squared $2$-Wasserstein distance between Gaussian approximations of feature distributions in the Inception representation \citep{heusel2017gans}. By contrast, small KL divergence alone does not, without further integrability assumptions, provide a uniform quantitative control of means, covariances, or Euclidean transportation error \citep{beyler2025convergence}. These observations motivate a direct $W_2$ analysis rather than obtaining geometric conclusions indirectly from a divergence bound.

\paragraph{Low-dimensional structure of data.}
Worst-case rates expressed in the ambient dimension can be highly pessimistic for modern data. 
The manifold hypothesis asserts that high-dimensional observations are often concentrated on, 
or close to, a set describable by substantially fewer effective variables. 
This effective number of variables is commonly called the intrinsic dimension. 
For example, although an ImageNet image 
 lies in an ambient space of dimension $224\times224\times3 = 150{,}528$, 
\citet{pope2021intrinsic} empirically estimate the intrinsic dimension of ImageNet to be approximately $43$.
Similar phenomena are observed for MNIST and CIFAR-10. This gap raises a fundamental question: can the convergence complexity of a diffusion sampler depend on the intrinsic dimension of the data rather than on the ambient dimension of the representation?

For KL divergence, recent work has answered this question
affirmatively. \citet{li2024adapting} studied adaptation to unknown
low-dimensional structure, \citet{azangulov2024manifold} analyzed
diffusion models under a manifold hypothesis, and
\citet{huang2026optimal} established a sharp intrinsic-dimension
dependence for the original DDPM sampler under a metric-entropy
assumption on the data support. In the latter framework, the support
need not be a smooth manifold and the target measure need not have a
Lebesgue density. The low-dimensional complexity is encoded by a
covering-number bound, and the KL bound depends on an intrinsic
quantity of order $k$, up to logarithmic factors, rather than on the
ambient dimension.

For total variation, \citet[Theorem~3]{li2025od} directly analyze the
DDPM recursion under a metric-entropy notion of intrinsic dimension.
With the coefficient design adapted to low-dimensional structure,
their exact-score bound has rate $\widetilde O(k/N)$.
\citet{liang2025lowdimtv} also develop a direct TV analysis under a
metric-entropy assumption, covering both the deterministic DDIM
sampler and the stochastic DDPM sampler. With exact scores, they show
that both samplers require $\widetilde O(k/\varepsilon)$ iterations
to achieve TV error $\varepsilon$, equivalently yielding the rate
$\widetilde O(k/N)$, and they also treat learned scores. These direct
TV rates improve on the order $\sqrt{k/N}$ that would follow by
applying Pinsker's inequality to an order-$k/N$ intrinsic-dimensional
KL bound.

A corresponding intrinsic-dimensional convergence theory
in 2-Wasserstein distance, however, was left open.

\paragraph{Our contribution.}
This paper establishes an intrinsic-dimension-dependent convergence
theory in $W_2$ for the original DDPM sampler. We use the
covering-number, or metric-entropy, notion of intrinsic dimension
adopted in the recent DDPM analyses of
\citet{li2024adapting}, \citet{huang2026optimal},
\citet{li2025od}, and \citet{liang2025lowdimtv}. The precise
formulations differ only in whether the covering bound is imposed at
a prescribed resolution or uniformly over resolutions: our one-scale
condition agrees with that of \citet{huang2026optimal}, closely
matches the one-scale formulations of \citet{li2024adapting} and
\citet{li2025od}, and is implied by the all-scale condition of
\citet{liang2025lowdimtv}. Thus, our Wasserstein result does not
require a stronger assumption on the low-dimensional structure than
these preceding KL and TV theories. Under this condition, $k$
represents the intrinsic complexity of the support. Unlike
smooth-manifold formulations such as \citet{azangulov2024manifold},
this framework allows the support to be highly irregular: the target
distribution need not be log-concave, possess a density, or lie on a
smooth manifold.

Our first main result gives a non-asymptotic $W_2$ guarantee after
early stopping, in which the time-discretization error is governed by
the intrinsic dimension, up to logarithmic factors. A corollary then
combines this sampling guarantee with the explicit early-stopping
error to control the distance from the original data distribution.
This separation identifies precisely where the ambient dimension
remains: in the Gaussian-initialization term of the post-early-stopping
bound and in the additional early-stopping term of the overall bound.

Our second main result applies the same global error accounting to KL
divergence. For the standard time mesh, it improves the sufficient
iteration complexity by one logarithmic factor relative to the
corresponding intrinsic-dimensional analysis of
\citet{huang2026optimal}, while preserving a nearly linear dependence
on $k$ up to logarithmic factors. Both results allow imperfect score
estimation. Together, they show that the original DDPM coefficients
can exploit unknown low-dimensional structure without requiring the
sampler to identify or parameterize that structure explicitly.

\paragraph{Organization.}
\Cref{sec:setup} introduces the forward diffusion, the reverse-time SDE, and the continuous-time interpolation of DDPM. \Cref{sec:assumptions} states the metric-entropy, mesh, and score assumptions. \Cref{sec:main} presents the Wasserstein and KL bounds. \Cref{sec:proof} proves the Wasserstein theorem by combining the reverse posterior martingale, a quadratic discrete recursion, and the information-theoretic discretization estimate. The appendices give the DDPM-interpolation, martingale, and recursion proofs, the metric-entropy information bound, the Stieltjes/I--MMSE analysis, the Girsanov proof of the KL theorem, and the proof of the overall Wasserstein corollary.

\section{Preliminaries}
\label{sec:setup}

\subsection{Forward Ornstein--Uhlenbeck process}

Let $X_0\sim\pdata$ be an $\R^d$-valued random vector satisfying $\E\left[\norm{X_0}^2\right]<\infty$. Consider the Ornstein--Uhlenbeck process
\begin{equation}
    \dd X_t=-X_t\dd t+\sqrt{2}\,\dd W_t,
    \qquad t\in[0,T],
    \label{eq:forward-sde}
\end{equation}
where $W$ is a standard $d$-dimensional Brownian motion independent of $X_0$. Its marginal law admits the representation
\begin{equation}
    X_t\overset{\mathrm d}=e^{-t}X_0+\sigma_t Z,
    \qquad
    \sigma_t:=\sqrt{1-e^{-2t}},
    \qquad Z\sim\N(0,\Id),
    \label{eq:forward-marginal}
\end{equation}
with $Z$ independent of $X_0$. For $t>0$, let $q_t$ be the density of $X_t$ and define the score
\begin{equation}
    s_t(x):=\nabla\log q_t(x).
    \label{eq:true-score}
\end{equation}

\subsection{Reverse-time diffusion}

The reverse-time SDE associated with \eqref{eq:forward-sde} is
\begin{equation}
    \dd Y_t=\bigl\{Y_t+2s_{T-t}(Y_t)\bigr\}\dd t+\sqrt{2}\,\dd B_t,
    \qquad t\in[0,T),
    \label{eq:reverse-sde}
\end{equation}
where $B$ is a $d$-dimensional Brownian motion. We first clarify the
well-posedness and the time-reversal interpretation of this equation. Fix
$\tau<T$. Since $q_r$ is a strictly positive Gaussian convolution for every
$r>0$, the map $(r,x)\mapsto s_r(x)=\nabla\log q_r(x)$ is continuously
differentiable in $x$ on $(0,\infty)\times\R^d$. Consequently, for every
$R<\infty$, there exists $K_{\tau,R}<\infty$ such that
\begin{align}
    \norm{
        \bigl(x+2s_{T-t}(x)\bigr)
        -
        \bigl(y+2s_{T-t}(y)\bigr)
    }
    &\le K_{\tau,R}\norm{x-y},
    \qquad
    0\le t\le\tau,\quad
    \norm{x},\norm{y}\le R.
    \notag
\end{align}
The dispersion coefficient $\sqrt{2}\,\Id$ is constant. Hence
\citet[Theorem~5.2.5]{karatzas1998brownian} implies pathwise
(strong) uniqueness for \eqref{eq:reverse-sde} on $[0,\tau]$. Since
$\tau<T$ is arbitrary, pathwise uniqueness holds on $[0,T)$.

Weak existence is supplied directly by time reversal. Indeed, the forward
SDE \eqref{eq:forward-sde} has drift $x\mapsto-x$ and the constant
dispersion coefficient $\sqrt{2}\,\Id$; these coefficients are globally
Lipschitz and have linear growth. Moreover, the Gaussian-smoothed
densities $q_r$, $r>0$, satisfy the local
regularity condition in \citet[Theorem~2.1]{haussmann1986time}. That theorem,
applied on every interval $[0,\tau]$ with $\tau<T$, shows that the canonical
time-reversed process $(X_{T-t})_{0\le t\le\tau}$ is again a
diffusion. Its dispersion coefficient is $\sqrt{2}\,\Id$, and its reverse
drift is
\begin{align}
    -(-x)
    +\frac{\nabla\cdot\bigl(2\Id q_{T-t}(x)\bigr)}
           {q_{T-t}(x)}
    &=
    x+2\nabla\log q_{T-t}(x)
    =
    x+2s_{T-t}(x).
    \notag
\end{align}
Thus, for a suitable Brownian motion $B$, the process
$(X_{T-t})_{0\le t\le\tau}$ is a weak solution of
\eqref{eq:reverse-sde}. Combining this weak existence with the
pathwise uniqueness established above,
\citet[Proposition~5.3.20]{karatzas1998brownian} implies uniqueness
in law for \eqref{eq:reverse-sde} on every $[0,\tau]$, $\tau<T$.
Consequently, every weak solution $Y$ of \eqref{eq:reverse-sde} with
$Y_0\sim\Law(X_T)$ has the same law on
$C([0,\tau];\R^d)$ as $(X_{T-t})_{0\le t\le\tau}$. In particular,
the distributional time-reversal identity
\begin{equation}
    Y_t\overset{\mathrm d}=X_{T-t},
    \qquad 0\le t<T,
    \label{eq:time-reversal-marginal}
\end{equation}
holds for the solution of \eqref{eq:reverse-sde}.

For later use, define the posterior mean and posterior covariance
\begin{equation}
    \mu_t(x):=\E\left[X_0\mid X_t=x\right],
    \qquad
    \Sigma_t(x):=\Cov(X_0\mid X_t=x).
    \label{eq:posterior-defs}
\end{equation}
Tweedie's formula \citep{efron2011tweedie} gives
\begin{equation}
    \mu_t(x)=\frac{x+\sigma_t^2s_t(x)}{e^{-t}},
    \qquad
    s_t(x)=\frac{e^{-t}}{\sigma_t^2}\mu_t(x)-\frac{x}{\sigma_t^2}.
    \label{eq:tweedie}
\end{equation}
Consequently, \eqref{eq:reverse-sde} can be rewritten as
\begin{equation}
    \dd Y_t=
    \left\{
       \left(1-\frac{2}{\sigma_{T-t}^2}\right)Y_t
       +\frac{2e^{-(T-t)}}{\sigma_{T-t}^2}\mu_{T-t}(Y_t)
    \right\}\dd t+\sqrt{2}\,\dd B_t.
    \label{eq:reverse-posterior}
\end{equation}

\subsection{DDPM and its semilinear interpolation}

Choose grid points
\begin{equation}
    0=t_0<t_1<\cdots<t_N<t_{N+1}=T.
    \label{eq:grid}
\end{equation}
Set
\begin{equation}
    \alpha_n:=e^{-2(t_{n+1}-t_n)},
    \qquad
    \bar\alpha_n:=\prod_{i=n}^{N}\alpha_i=e^{-2(T-t_n)}.
    \label{eq:alpha}
\end{equation}
We first define the discrete DDPM sampler. Given score estimators
$\hat s_t$, let $(\widetilde Y_n)_{n=0}^N$ be the discrete-time
sequence generated by
\begin{align}
    \widetilde Y_0&\sim\N(0,\Id),\notag\\
    \widetilde Y_{n+1}
    &=\frac{1}{\sqrt{\alpha_n}}
       \left\{\widetilde Y_n+(1-\alpha_n)
       \hat s_{T-t_n}(\widetilde Y_n)\right\}
       +\sqrt{\frac{(1-\alpha_n)(1-\bar\alpha_{n+1})}
       {1-\bar\alpha_n}}\,Z_n,
\label{eq:ddpm-update}
\end{align}
for $n=0,\ldots,N-1$, where $Z_0,\dots,Z_{N-1}$ are independent
standard Gaussian vectors and are independent of $\widetilde Y_0$.
This is the DDPM reverse sampling update introduced by
\citet{ho2020ddpm}, written in the Ornstein--Uhlenbeck time
parameterization of \eqref{eq:alpha} and in terms of the learned
score $\hat s_{T-t_n}$.
Here the subscript $n$ is an iteration index; at this point no
continuous-time process is being introduced.

Define the posterior-mean estimator induced by $\hat s_t$ as
\begin{equation}
    \hat\mu_t(x):=\frac{x+\sigma_t^2\hat s_t(x)}{e^{-t}}.
    \label{eq:muhat}
\end{equation}
We next introduce a separate continuous-time process $\hat Y$ whose
grid-point values will be shown to reproduce the discrete recursion
\eqref{eq:ddpm-update}. Let $\hat Y_{t_0}\sim\N(0,\Id)$ be independent
of a $d$-dimensional Brownian motion $B$.
SDE~\eqref{eq:reverse-posterior} separates the exact reverse
drift into a time-dependent linear term and a nonlinear
posterior-mean term. To define the continuous-time interpolation, we keep the
linear coefficient and the diffusion coefficient exactly as in
\eqref{eq:reverse-posterior}, replace $Y_t$ by $\hat Y_t$, and freeze
only the learned posterior-mean term at the left endpoint:
\[
    \mu_{T-t}(Y_t)
    \quad\longrightarrow\quad
    \hat\mu_{T-t_n}(\hat Y_{t_n}),
    \qquad t\in[t_n,t_{n+1}].
\]
This gives the following semilinear SDE on $[t_n,t_{n+1}]$:
\begin{equation}
    \dd\hat Y_t=
    \left\{
       \left(1-\frac{2}{\sigma_{T-t}^2}\right)\hat Y_t
       +\frac{2e^{-(T-t)}}{\sigma_{T-t}^2}
         \hat\mu_{T-t_n}(\hat Y_{t_n})
    \right\}\dd t+\sqrt{2}\,\dd B_t.
    \label{eq:interpolation}
\end{equation}
Thus, \eqref{eq:interpolation} has the same semilinear structure as
\eqref{eq:reverse-posterior}: the linear part is integrated exactly,
whereas the nonlinear learned posterior mean is held constant on each
grid interval. The following proposition establishes the connection
between this SDE and the discrete algorithm \eqref{eq:ddpm-update}.
The grid-point equivalence underlying this result is also proved in
\citet{huang2026optimal}. For completeness,
\Cref{app:interpolation-proof} gives a proof that makes explicit both
the construction of the Gaussian vectors from the Brownian increments
and the almost-sure coupling with the discrete recursion.

\begin{proposition}[Exact grid-point realization of DDPM]
\label{prop:interpolation}
Construct $\hat Y$ recursively from \eqref{eq:interpolation}. Then the
Brownian increments on the grid intervals determine independent
random vectors $Z_0,\ldots,Z_{N-1}\sim\N(0,\Id)$, independent of
$\hat Y_{t_0}$, such that, for every $n=0,\ldots,N-1$, the
grid-point values of the SDE satisfy
\begin{align}
    \hat Y_{t_{n+1}}
    ={}&\frac{1}{\sqrt{\alpha_n}}
       \left\{\hat Y_{t_n}+(1-\alpha_n)
       \hat s_{T-t_n}(\hat Y_{t_n})\right\}\notag\\
    &+\sqrt{\frac{(1-\alpha_n)(1-\bar\alpha_{n+1})}
       {1-\bar\alpha_n}}\,Z_n.\notag
\end{align}
Consequently, the initial conditions and Gaussian innovations of the
two recursions have the same joint distribution, and hence
\begin{align}
    (\widetilde Y_0,\ldots,\widetilde Y_N)
    \overset{\mathrm d}{=}
    (\hat Y_{t_0},\ldots,\hat Y_{t_N}).
    \notag
\end{align}
Thus, \eqref{eq:interpolation} is a continuous-time interpolation of
the original DDPM sampler in distribution.
\end{proposition}

\section{Intrinsic dimension and analytic assumptions}
\label{sec:assumptions}

\subsection{Metric-entropy structure}

The ambient dimension $d$ records the number of coordinates, but it
need not reflect the number of geometrically distinct configurations
occupied by the data. To describe low-dimensional structure without
assuming that the support is a smooth manifold or even that
$\pdata$ has a density, we instead measure the complexity of the
support at a prescribed Euclidean resolution. Following the standard
terminology in \citet[Chapter~5]{wainwright2019high} and the
low-dimensional DDPM formulation of
\citet[Section~4.1.1]{huang2026optimal}, we use the following notion.

\begin{definition}[Covering number and metric entropy]
For $A\subset\R^d$ and $\varepsilon>0$, define
\[
    N_{\mathrm{cov}}(A,\norm{\cdot},\varepsilon)
\]
to be the smallest integer $m$ for which there exist
$x_1,\dots,x_m\in\R^d$ such that
\[
    A\subset\bigcup_{i=1}^m B(x_i,\varepsilon).
\]
Here,
\[
    B(x_i,\varepsilon)
    :=\left\{x\in\R^d:\norm{x-x_i}\le\varepsilon\right\}.
\]
The corresponding Euclidean metric entropy of $A$ at resolution
$\varepsilon$ is
\[
    \log N_{\mathrm{cov}}(A,\norm{\cdot},\varepsilon).
\]
\end{definition}

Let $\mathcal X:=\supp(\pdata)$.
The covering number counts how many representatives are needed to
describe the set up to accuracy $\varepsilon$, while its logarithm
measures the logarithmic description length of this finite
representation. This is the scale-sensitive notion of complexity
needed in our analysis. Indeed, an $\varepsilon_0$-net allows each
possible value of $X_0$ to be assigned to one of at most
$N_{\mathrm{cov}}(\mathcal X,\norm{\cdot},\varepsilon_0)$
representatives. In the information-theoretic discretization argument,
the logarithm of this number controls the information carried by the
representative index, while the distance from $X_0$ to its
representative controls the remaining conditional uncertainty. Metric
entropy therefore provides the bridge from the geometry of the data
support to the intrinsic-dimensional $k\log k$ error bound.

\begin{assumption}[Metric-entropy dimension]
\label{ass:entropy}
Let $k\ge2$. There are fixed constants $C_0,C_{\mathrm{cov}}>0$ such that, with
$\varepsilon_0=k^{-C_0}$,
\begin{equation}
    \log N_{\mathrm{cov}}(\mathcal X,\norm{\cdot},\varepsilon_0)
    \le C_{\mathrm{cov}}k\log(1/\varepsilon_0).
    \label{eq:entropy-assumption}
\end{equation}
\end{assumption}

\begin{assumption}[Bounded support]
\label{ass:support}
For a fixed constant $C_R>0$,
\begin{equation}
    \sup_{x\in\mathcal X}\norm{x}\le R,
    \qquad R\le k^{C_R}.
    \label{eq:support-bound}
\end{equation}
\end{assumption}

Assumption~\ref{ass:entropy} means that, at resolution
$\varepsilon_0$, the support can be represented by a net containing
at most
\[
    \exp\left\{
        C_{\mathrm{cov}}k\log(1/\varepsilon_0)
    \right\}
    =k^{C_{\mathrm{cov}}C_0k}
\]
points. Thus, $k$ measures the logarithmic complexity of the support
at the resolution relevant to our analysis, rather than requiring the
support to have topological dimension $k$. The condition is imposed
only at the single scale $\varepsilon_0$ and therefore does not demand
regular behavior at every finer scale. Assumption~\ref{ass:support}
is also permissive: because $C_R$ is any fixed constant, the support
radius may be an arbitrarily large fixed polynomial in $k$.

As discussed in \citet[Section~4.1.1]{huang2026optimal}, this
metric-entropy framework covers a variety of familiar
low-dimensional structures:
\begin{itemize}[leftmargin=2em]
    \item A bounded subset of a $k$-dimensional linear subspace has
    metric entropy of order $k\log(1/\varepsilon_0)$, up to the
    contribution of its polynomial support radius.

    \item A sufficiently regular compact $k$-dimensional manifold
    satisfies the condition when its volume and inverse reach are
    polynomially controlled. Small neighborhoods of such a manifold
    can also be accommodated, so the data need not lie exactly on the
    manifold.

    \item More general sets described by doubling dimension or
    Minkowski dimension $k$ also satisfy the same type of
    covering-number bound.

    \item For $s$-sparse vectors in $\R^d$, the effective
    metric-entropy dimension is of order $s\log d$. The condition is
    also stable under unions of up to $k^{O(k)}$ component sets of
    comparable metric entropy.
\end{itemize}
Consequently, no manifold parametrization, smoothness, connectedness,
or absolute continuity of $\pdata$ is required; the support may be
singular, disconnected, and geometrically irregular.

\subsection{Time mesh and score assumptions}

Recall the reverse-time grid \eqref{eq:grid}. The exact data endpoint
corresponds to $t=T$, where the remaining forward noising time
$T-t$ vanishes. Near this endpoint, the reverse coefficients vary
rapidly and may become singular when the data distribution has no
density. We therefore stop the reverse sampler at
$t_N=T-\delta$ and call
\begin{equation}
    \delta:=T-t_N.
    \label{eq:delta}
\end{equation}
the early-stopping time. Thus, the main theorem compares the sampler
with the positively smoothed marginal $\Law(X_\delta)$, while
\Cref{cor:overall-W2} separately adds the error between $\pdata$ and
$\Law(X_\delta)$.

Following the time discretization used by
\citet{benton2024nearly}, define the relative mesh size $\kappa$ as
the smallest number for which
\begin{equation}
    t_{n+1}-t_n\le\kappa\min\{1,T-t_n\},
    \qquad 0\le n\le N-1.
    \label{eq:mesh-assumption}
\end{equation}
This condition combines two different resolutions. As long as
$T-t_n\ge1$, it requires the ordinary time increment
$t_{n+1}-t_n$ to be at most $\kappa$. Once $T-t_n<1$, it instead
controls the increment relative to the remaining forward time. In
particular,
\[
    T-t_{n+1}\ge(1-\kappa)(T-t_n)
\]
in this second regime. Hence the grid becomes progressively finer near
the data endpoint, precisely where the reverse drift and the
signal-to-noise ratio change most rapidly. This relative control is the
mesh property used later to accumulate the local discretization errors
without paying an additional factor proportional to the number of
intervals.

For orientation, consider the standard two-stage mesh used in recent
diffusion-model analyses. Assuming that $T>1$ and $N$ is even, it is
given by
\begin{equation}
    t_n=
    \begin{cases}
        \displaystyle \frac{2(T-1)n}{N},
        &0\le n\le N/2,\\[6pt]
        \displaystyle T-\delta^{\,2(n-N/2)/N},
        &N/2<n\le N.
    \end{cases}
    \label{eq:two-stage-mesh}
\end{equation}
The first half is uniformly spaced on $[0,T-1]$, whereas in the second
half the remaining time $T-t_n$ decreases geometrically from $1$ to
$\delta$. As shown by
\citet[Corollary~1]{benton2024nearly}, this mesh satisfies
\begin{equation}
    \kappa\asymp
    \frac{T+\log(1/\delta)}{N}.
    \label{eq:two-stage-kappa}
\end{equation}
Thus, $\kappa$ is a dimensionless measure of the effective mesh
resolution: it is of order $N^{-1}$ when $T$ and $\delta$ are fixed,
while \eqref{eq:two-stage-kappa} records the additional refinement
needed when $T$ grows or $\delta$ decreases.

It is useful first to relate the score assumptions to how the score is
learned. At a fixed forward time $r$, the population denoising
score-matching objective is equivalent, up to a term independent of
the learned score, to the squared error
\citep{song2021score}
\begin{align}
    \E\left[
        \norm{s_r(X_r)-\hat s_r(X_r)}^2
    \right].
    \notag
\end{align}
Thus, Assumption~\ref{ass:score-error} evaluates the learned score
under the same forward marginal that appears in the score-matching
loss. By contrast, an error bound under the approximate sampling
distribution does not directly correspond to this training objective:
the law of $\hat Y_{t_n}$ is generated by the learned reverse dynamics
and is not the forward noising distribution used for training.

Existing Wasserstein analyses control the learned score in one of
these two ways. The first is nevertheless to evaluate its error
directly under the approximate sampling distribution. In the notation
of this paper, this amounts to assuming, for every
$n=0,\ldots,N-1$, a bound of the form
\begin{align}
    \E\left[
        \norm{
            s_{T-t_n}(\hat Y_{t_n})
            -\hat s_{T-t_n}(\hat Y_{t_n})
        }^2
    \right]
    \le \varepsilon_n^2.
    \notag
\end{align}
Here the expectation is taken at the actual inputs encountered by the
approximate sampler, whose distribution itself depends on the learned
score. The second approach evaluates the score error under the forward
marginal $\Law(X_{T-t_n})$, as in
Assumption~\ref{ass:score-error}. Because the score is then evaluated
under a distribution different from the approximate sampling
distribution, existing Wasserstein analyses pair this
training-aligned error bound with a spatial regularity condition on
the learned score, such as Assumption~\ref{ass:lipschitz}, in order to
compare the two sampling laws. This distinction is also discussed by
\citet[Section~2.2]{koike2026wasserstein}. Thus, the preceding
Wasserstein literature assumes either the displayed score-error bound
under the approximate sampling distribution, or the combination of
Assumptions~\ref{ass:lipschitz} and \ref{ass:score-error}. We adopt the
latter formulation.

\begin{assumption}[Spatial regularity of the learned score]
\label{ass:lipschitz}
There is a measurable function $L:[\delta,T]\to[0,\infty)$ such that
\begin{equation}
    \norm{\hat s_r(x)-\hat s_r(y)}
    \le L(r)\norm{x-y},
    \qquad r\in[\delta,T],\quad x,y\in\R^d.
    \label{eq:lipschitz-score}
\end{equation}
\end{assumption}

\begin{assumption}[Score estimation along the forward marginals]
\label{ass:score-error}
For $n=0,\dots,N-1$, there are finite numbers $\varepsilon_n\ge0$ such that
\begin{equation}
    \E\left[
    \norm{s_{T-t_n}(X_{T-t_n})-
    \hat s_{T-t_n}(X_{T-t_n})}^2
    \right]\le\varepsilon_n^2.
    \label{eq:score-error}
\end{equation}
\end{assumption}

\section{Main result}
\label{sec:main}

For $u>0$, define the Ornstein--Uhlenbeck signal-to-noise ratio
\begin{equation}
    \gamma(u):=\frac{1-\sigma_u^2}{\sigma_u^2}
    =\frac{1}{e^{2u}-1}.
    \label{eq:gamma}
\end{equation}

\begin{theorem}[Post-early-stopping Wasserstein bound]
\label{thm:main}
Suppose Assumptions~\ref{ass:entropy}--\ref{ass:score-error} hold. Assume $T>1$, $k\ge2$, and that for fixed constants $C_d,C_\delta>0$,
\begin{equation}
    d\le k^{C_d},
    \qquad
    k^{-C_\delta}\le\delta<1,
    \qquad
    \kappa\le0.9,
    \qquad
    2C_0>C_\delta.
    \label{eq:poly-regime}
\end{equation}
Then there is a constant $C>0$, depending only on the fixed constants in the assumptions, such that
\begin{align}
\Wtwo\bigl(\Law(X_\delta),\Law(\widetilde Y_N)\bigr)
\le{}&
\exp\left\{t_N+2\sum_{i=0}^{N-1}(t_{i+1}-t_i)L(T-t_i)\right\}
\Biggl[
\underbrace{\left\{e^{-2T}\E\left[\norm{X_0}^2\right]+(1-\sigma_T)^2d\right\}^{1/2}}_{\text{Gaussian-initialization error}}
\notag\\
&
+\underbrace{\sum_{j=0}^{N-1}e^{t_{j+1}-t_j}
\left\{1-e^{-2(t_{j+1}-t_j)}\right\}\varepsilon_j}_{\text{score-estimation error}}
+\underbrace{
 C\left(\kappa k\log k\right)^{1/2}
}_{\text{time-discretization error}}
\quad \Biggr].
\label{eq:post-early-stopping-bound}
\end{align}
\end{theorem}

\begin{corollary}[Overall Wasserstein bound]
\label{cor:overall-W2}
Under the assumptions of \Cref{thm:main},
\begin{align}
\Wtwo\bigl(\pdata,\Law(\widetilde Y_N)\bigr)
\le
&\exp\left\{t_N+2\sum_{i=0}^{N-1}(t_{i+1}-t_i)L(T-t_i)\right\}
\Biggl[
\underbrace{\left\{e^{-2T}\E\left[\norm{X_0}^2\right]+(1-\sigma_T)^2d\right\}^{1/2}}_{\text{Gaussian-initialization error}}
\notag\\
&
+\underbrace{\sum_{j=0}^{N-1}e^{t_{j+1}-t_j}
\left\{1-e^{-2(t_{j+1}-t_j)}\right\}\varepsilon_j}_{\text{score-estimation error}}
+\underbrace{
 C\left(\kappa k\log k\right)^{1/2}
}_{\text{time-discretization error}}
\quad \Biggr]\notag\\
&+\underbrace{\left\{(1-e^{-\delta})^2
\E\left[\norm{X_0}^2\right]+\sigma_\delta^2d\right\}^{1/2}
}_{\text{early-stopping error}}
\label{eq:overall-W2-bound}
\end{align}
\end{corollary}

The proof is deferred to \Cref{app:overall-W2}.

\begin{theorem}[Improved intrinsic-dimension-dependent KL bound]
\label{thm:main-KL}
Suppose Assumptions~\ref{ass:entropy}, \ref{ass:support}, and
\ref{ass:score-error} hold, and assume \eqref{eq:poly-regime}. Let
\[
    p_{\mathrm{out}}
    :=\Law(\widetilde Y_N)
    =\Law(\hat Y_{t_N}),
\]
where the equality follows from Proposition~\ref{prop:interpolation}.
Then
\begin{align}
    \KL(q_\delta\Vert p_{\mathrm{out}})
    \le C\Bigl\{
       &\kappa k\log k+{\sum_{j=0}^{N-1}e^{t_{j+1}-t_j}\left\{1-e^{-2(t_{j+1}-t_j)}\right\}\varepsilon_j}+\left(d+\E\left[\norm{X_0}^2\right]\right)e^{-2T}
    \Bigr\}.
    \label{eq:main-KL}
\end{align}
The constant $C$ depends only on the fixed constants in the
assumptions.

For the standard two-stage mesh, for which
\[
    \kappa\asymp\frac{T+\log(1/\delta)}{N},
\]
let $\varepsilon\in(0,1]$, suppose
$\scoreerr^2\lesssim\varepsilon^2$, and choose
\[
    T\gtrsim
    \log\left(
       \frac{d+\E\left[\norm{X_0}^2\right]+1}{\varepsilon}
    \right).
\]
Then
\begin{equation}
    N\gtrsim
    \frac{k\log k\,\{T+\log(1/\delta)\}}{\varepsilon^2}
    \label{eq:KL-iteration-general}
\end{equation}
implies
$\KL(q_\delta\Vert p_{\mathrm{out}})\lesssim\varepsilon^2$.
In the polynomial regime \eqref{eq:poly-regime}, this is ensured by
\begin{equation}
    N\gtrsim
    \frac{k\log k\,\log(k/\varepsilon)}{\varepsilon^2}.
    \label{eq:KL-iteration-sharp}
\end{equation}
Since $\log k\le\log(k/\varepsilon)$ for $\varepsilon\le1$, the
simpler sufficient condition is
\begin{equation}
    N\gtrsim
    \frac{k\log^2(k/\varepsilon)}{\varepsilon^2}.
    \label{eq:KL-iteration}
\end{equation}
\end{theorem}

The proof of \Cref{thm:main-KL} is given in \Cref{app:KL}. Compared
with the earlier interval-by-interval covariance estimate, whose
standard-mesh specialization yields the sufficient order
$k\log^3(k/\varepsilon)/\varepsilon^2$, \eqref{eq:KL-iteration}
removes one logarithmic factor.

\begin{remark}[Where the ambient dimension remains]
The squared discretization term in
\eqref{eq:post-early-stopping-bound} is governed by
$\kappa k\log k$, obtained from the mutual-information estimate.
Within the main theorem, the ambient dimension $d$ remains only in
the Gaussian-initialization error. The overall bound
\eqref{eq:overall-W2-bound} contains in addition the early-stopping
error. These two terms arise before the intrinsic-dimensional
posterior fluctuation estimate is used. Choosing $T$ large and
$\delta$ small controls them, but the required choices may depend on
$d$.
\end{remark}

\begin{remark}[Score regularity]
The Lipschitz assumption is imposed only on the learned score, not on the target density or the true score. It enters through the pathwise stability of the synchronous coupling. Removing this assumption would require a different stability mechanism, such as one-sided contractivity or a localized argument.
\end{remark}

\section{Proof of the Wasserstein theorem}
\label{sec:proof}

We work on the probability space carrying the forward
Ornstein--Uhlenbeck process $(X_t)_{0\le t\le T}$ and define its
canonical time reversal by
\[
    Y_t:=X_{T-t},
    \qquad 0\le t\le t_N.
\]
Thus,
\[
    Y_t=X_{T-t}
    \qquad\text{almost surely},
\]
and, in particular,
\[
    Y_{t_N}=X_\delta
    \qquad\text{almost surely}.
\]

By the time-reversal theorem
\citep{anderson1982reverse,haussmann1986time}, there exists a
$d$-dimensional Brownian motion $B$, with respect to the completed
natural reverse filtration
\[
    \mathcal G_t:=\sigma\{Y_s:0\le s\le t\},
\]
such that $Y$ satisfies \eqref{eq:reverse-posterior}. Notice that
$Y$ is not chosen as an arbitrary weak solution of the reverse SDE:
it is fixed as the canonical time reversal of the forward process.
Consequently, no uniqueness in law for the reverse SDE is needed in
order to identify its marginals with those of the forward process.

Because $Y_{t_N}=X_\delta$ almost surely, the comparison required for
\Cref{thm:main} is exactly the comparison between the endpoint laws of
the exact and approximate reverse processes. We next construct a
synchronous coupling of these two processes.

\subsection{Synchronous coupling}

Let
\[
    \gamma_d:=\N(0,\Id).
\]
Since $T>0$, the forward marginal $q_T=\Law(X_T)$ is the convolution
of the law of $e^{-T}X_0$ with the nondegenerate Gaussian law
$\N(0,\sigma_T^2\Id)$. Hence $q_T$ is absolutely continuous with
respect to Lebesgue measure. Moreover, both $q_T$ and $\gamma_d$
have finite second moments.

Therefore, by Brenier's theorem
\citep[Theorem~9.4]{villani2009optimal}, there exists a Borel
measurable optimal transport map
\[
    \mathcal T_T:\R^d\to\R^d
\]
such that
\[
    (\mathcal T_T)_{\#}q_T=\gamma_d
\]
and
\[
    W_2^2(q_T,\gamma_d)
    =
    \int_{\R^d}
    \norm{y-\mathcal T_T(y)}^2\,q_T(\dd y).
\]

We define the initial value of the approximate reverse process by
\begin{equation}
    \hat Y_0:=\mathcal T_T(Y_0).
    \label{eq:explicit-initial-coupling}
\end{equation}
Then
\[
    \hat Y_0\sim\gamma_d=\N(0,\Id)
\]
and
\begin{equation}
    \E\left[\norm{Y_0-\hat Y_0}^2\right]
    =
    W_2^2(q_T,\gamma_d).
    \label{eq:optimal-initial-cost}
\end{equation}
In particular, $\hat Y_0$ is $\mathcal G_0$-measurable.

Starting from \eqref{eq:explicit-initial-coupling}, construct
$\hat Y$ recursively on each interval $[t_n,t_{n+1}]$ using
\eqref{eq:interpolation}, driven by the same Brownian motion $B$ as
the canonical reverse process $Y$. On $[t_n,t_{n+1}]$, the random
vector
\[
    \hat\mu_{T-t_n}(\hat Y_{t_n})
\]
is frozen at the left endpoint. Conditional on $\mathcal G_{t_n}$,
\eqref{eq:interpolation} is therefore a linear SDE and is uniquely
defined by its variation-of-constants formula.

Since $B$ is a Brownian motion with respect to
$(\mathcal G_t)_{0\le t\le t_N}$, its increments are independent of
$\mathcal G_0$ and hence independent of $\hat Y_0$. It follows from
Proposition~\ref{prop:interpolation} that the grid-point values
$(\hat Y_{t_n})_{n=0}^N$ have the same distribution as the DDPM
sampler \eqref{eq:ddpm-update}.

The resulting synchronous coupling gives
\begin{equation}
    \Wtwo\bigl(\Law(Y_{t_N}),\Law(\hat Y_{t_N})\bigr)
    \le
    \left(
        \E\left[\norm{Y_{t_N}-\hat Y_{t_N}}^2\right]
    \right)^{1/2}.
    \label{eq:w2-coupling}
\end{equation}

We finally record the filtration associated with this construction.
Since $\hat Y_0$ is $\mathcal G_0$-measurable and $\hat Y$ is
constructed recursively from $\hat Y_0$ and the Brownian path
$(B_s)_{0\le s\le t}$, we have
\[
    \sigma\{\hat Y_s:0\le s\le t\}
    \subseteq\mathcal G_t.
\]
Consequently, if
\[
    \mathcal F_t
    :=
    \sigma\{Y_s,\hat Y_s:0\le s\le t\},
\]
with the usual completion, then
\[
    \mathcal F_t=\mathcal G_t.
\]
Thus, adjoining the approximate process $\hat Y$ does not reveal
any information beyond the path of the canonical reverse process
$Y$.

\subsection{Posterior-mean martingale and one-step error decomposition}

\begin{lemma}[Reverse posterior martingale]
\label{lem:martingale}
The process
\[
    \bigl(\mu_{T-t}(Y_t)\bigr)_{0\le t\le t_N}
\]
is a square-integrable martingale. Moreover, for $0\le s\le t\le t_N$,
\begin{equation}
    \E\left[
        \norm{\mu_{T-t}(Y_t)-\mu_{T-s}(Y_s)}^2
    \right]
    =
    \E\left[\Tr\Sigma_{T-s}(Y_s)\right]
    -
    \E\left[\Tr\Sigma_{T-t}(Y_t)\right].
    \label{eq:martingale-isometry}
\end{equation}
\end{lemma}

A closely related version of Lemma~\ref{lem:martingale} appears in
\citet{huang2026optimal}. For completeness, and to make the argument
self-contained under the synchronous coupling used here, we provide
a proof in \Cref{app:martingale}.

We now use this martingale structure to decompose the one-step error.
Applying the variation-of-constants formula to the posterior form
\eqref{eq:reverse-posterior} of the exact reverse SDE on the grid
interval $[t_n,t_{n+1}]$ from \eqref{eq:grid} gives
\begin{align}
Y_{t_{n+1}}
={}&e^{-(t_{n+1}-t_n)}
\frac{\sigma_{T-t_{n+1}}^2}{\sigma_{T-t_n}^2}Y_{t_n}
+\int_{t_n}^{t_{n+1}}
 e^{-(t_{n+1}-u)}
 \frac{\sigma_{T-t_{n+1}}^2}{\sigma_{T-u}^2}
 \frac{2e^{-(T-u)}}{\sigma_{T-u}^2}\mu_{T-u}(Y_u)\dd u\notag\\
&+\sqrt{2}\int_{t_n}^{t_{n+1}}
 e^{-(t_{n+1}-u)}
 \frac{\sigma_{T-t_{n+1}}^2}{\sigma_{T-u}^2}\dd B_u.
\label{eq:exact-mild}
\end{align}
For comparison, we recall the representation
\eqref{eq:interpolation-solution} of the approximate process. This
representation is obtained from the interpolation
\eqref{eq:interpolation}, whose frozen drift uses the learned
posterior mean \eqref{eq:muhat}:
\begin{align}
\hat Y_{t_{n+1}}
={}&e^{-(t_{n+1}-t_n)}
\frac{\sigma_{T-t_{n+1}}^2}{\sigma_{T-t_n}^2}
\hat Y_{t_n}
+\frac{
    e^{-(T-t_{n+1})}\{1-e^{-2(t_{n+1}-t_n)}\}
}{
    \sigma_{T-t_n}^2
}
\hat\mu_{T-t_n}(\hat Y_{t_n})\notag\\
&+\sqrt{2}\int_{t_n}^{t_{n+1}}
 e^{-(t_{n+1}-u)}
 \frac{\sigma_{T-t_{n+1}}^2}{\sigma_{T-u}^2}\dd B_u.
\label{eq:interpolation-solution}
\end{align}
Because the two processes are driven by the same Brownian motion,
the stochastic integrals cancel when these two representations are
subtracted. Hence
\begin{align}
Y_{t_{n+1}}-\hat Y_{t_{n+1}}
={}&e^{-(t_{n+1}-t_n)}
\frac{\sigma_{T-t_{n+1}}^2}{\sigma_{T-t_n}^2}
\bigl(Y_{t_n}-\hat Y_{t_n}\bigr)\notag\\
&+\int_{t_n}^{t_{n+1}}
 e^{-(t_{n+1}-u)}
 \frac{\sigma_{T-t_{n+1}}^2}{\sigma_{T-u}^2}
 \frac{2e^{-(T-u)}}{\sigma_{T-u}^2}
 \mu_{T-u}(Y_u)\dd u\notag\\
&-\frac{
    e^{-(T-t_{n+1})}\{1-e^{-2(t_{n+1}-t_n)}\}
}{
    \sigma_{T-t_n}^2
}
\hat\mu_{T-t_n}(\hat Y_{t_n}).\notag
\end{align}
We add and subtract $\mu_{T-t_n}(Y_{t_n})$ inside the drift
integral. This separates the frozen posterior-mean terms from their
temporal fluctuation:
\begin{align}
Y_{t_{n+1}}-\hat Y_{t_{n+1}}
={}&e^{-(t_{n+1}-t_n)}
\frac{\sigma_{T-t_{n+1}}^2}{\sigma_{T-t_n}^2}
\bigl(Y_{t_n}-\hat Y_{t_n}\bigr)
-\frac{
    e^{-(T-t_{n+1})}\{1-e^{-2(t_{n+1}-t_n)}\}
}{
    \sigma_{T-t_n}^2
}
\hat\mu_{T-t_n}(\hat Y_{t_n})\notag\\
&+\left\{
\int_{t_n}^{t_{n+1}}
 e^{-(t_{n+1}-u)}
 \frac{\sigma_{T-t_{n+1}}^2}{\sigma_{T-u}^2}
 \frac{2e^{-(T-u)}}{\sigma_{T-u}^2}\dd u
\right\}
\times
\mu_{T-t_n}(Y_{t_n})\notag\\
&+\int_{t_n}^{t_{n+1}}
 e^{-(t_{n+1}-u)}
 \frac{\sigma_{T-t_{n+1}}^2}{\sigma_{T-u}^2}
 \frac{2e^{-(T-u)}}{\sigma_{T-u}^2}
\times
\bigl\{
    \mu_{T-u}(Y_u)-\mu_{T-t_n}(Y_{t_n})
\bigr\}\dd u.\notag
\end{align}
The deterministic coefficient in the second term is
\begin{equation}
\int_{t_n}^{t_{n+1}}
 e^{-(t_{n+1}-u)}
 \frac{\sigma_{T-t_{n+1}}^2}{\sigma_{T-u}^2}
 \frac{2e^{-(T-u)}}{\sigma_{T-u}^2}\dd u
 =\frac{e^{-(T-t_{n+1})}\{1-e^{-2(t_{n+1}-t_n)}\}}
 {\sigma_{T-t_n}^2}.
\label{eq:kernel-integral}
\end{equation}
We denote the final integral above by the local remainder
\begin{align}
    R_n:=\int_{t_n}^{t_{n+1}}
 e^{-(t_{n+1}-u)}
 \frac{\sigma_{T-t_{n+1}}^2}{\sigma_{T-u}^2}
 \frac{2e^{-(T-u)}}{\sigma_{T-u}^2}
 \times\bigl\{\mu_{T-u}(Y_u)-\mu_{T-t_n}(Y_{t_n})\bigr\}\dd u.
\label{eq:Rn}
\end{align}
Equations~\eqref{eq:kernel-integral} and~\eqref{eq:Rn} therefore give
\begin{align}
Y_{t_{n+1}}-\hat Y_{t_{n+1}}
={}&e^{-(t_{n+1}-t_n)}
\frac{\sigma_{T-t_{n+1}}^2}{\sigma_{T-t_n}^2}
\bigl(Y_{t_n}-\hat Y_{t_n}\bigr)\notag\\
&+\frac{
    e^{-(T-t_{n+1})}\{1-e^{-2(t_{n+1}-t_n)}\}
}{
    \sigma_{T-t_n}^2
}
\bigl\{
    \mu_{T-t_n}(Y_{t_n})
    -\hat\mu_{T-t_n}(\hat Y_{t_n})
\bigr\}
+R_n.\notag
\end{align}
We next rewrite the posterior-mean difference in terms of the true
and learned scores. Applying the first identity in
\eqref{eq:tweedie} with
$t=T-t_n$ and $x=Y_{t_n}$ gives
\[
\mu_{T-t_n}(Y_{t_n})
=e^{T-t_n}
\left\{
    Y_{t_n}
    +\sigma_{T-t_n}^2s_{T-t_n}(Y_{t_n})
\right\}.
\]
Similarly, applying the definition \eqref{eq:muhat} with
$t=T-t_n$ and $x=\hat Y_{t_n}$ gives
\[
\hat\mu_{T-t_n}(\hat Y_{t_n})
=e^{T-t_n}
\left\{
    \hat Y_{t_n}
    +\sigma_{T-t_n}^2
        \hat s_{T-t_n}(\hat Y_{t_n})
\right\}.
\]
Subtracting these two identities yields
\begin{align}
\mu_{T-t_n}(Y_{t_n})
-\hat\mu_{T-t_n}(\hat Y_{t_n})
=
e^{T-t_n}\bigl(Y_{t_n}-\hat Y_{t_n}\bigr)
+e^{T-t_n}\sigma_{T-t_n}^2
\bigl\{
    s_{T-t_n}(Y_{t_n})
    -\hat s_{T-t_n}(\hat Y_{t_n})
\bigr\}.\notag
\end{align}
By the definition of the grid in \eqref{eq:grid},
\[
    e^{-(T-t_{n+1})}e^{T-t_n}=e^{t_{n+1}-t_n}.
\]
Therefore, multiplication by the coefficient in
\eqref{eq:kernel-integral} gives
\begin{align}
&\frac{
    e^{-(T-t_{n+1})}\{1-e^{-2(t_{n+1}-t_n)}\}
}{
    \sigma_{T-t_n}^2
}
\bigl\{
    \mu_{T-t_n}(Y_{t_n})
    -\hat\mu_{T-t_n}(\hat Y_{t_n})
\bigr\}\notag\\
&\quad=
\frac{e^{t_{n+1}-t_n}\{1-e^{-2(t_{n+1}-t_n)}\}}
{\sigma_{T-t_n}^2}
\bigl(Y_{t_n}-\hat Y_{t_n}\bigr)\notag\\
&\qquad\quad
{}+e^{t_{n+1}-t_n}\{1-e^{-2(t_{n+1}-t_n)}\}
\bigl\{
    s_{T-t_n}(Y_{t_n})
    -\hat s_{T-t_n}(\hat Y_{t_n})
\bigr\}.\notag
\end{align}
The first term on the right-hand side combines with the linear term
in the preceding display. Indeed, using
$\sigma_r^2=1-e^{-2r}$ from \eqref{eq:forward-marginal} and
$T-t_{n+1}=T-t_n-(t_{n+1}-t_n)$ from \eqref{eq:grid}, we obtain
\begin{align}
&e^{-(t_{n+1}-t_n)}
\frac{\sigma_{T-t_{n+1}}^2}{\sigma_{T-t_n}^2}
{}+\frac{e^{t_{n+1}-t_n}\{1-e^{-2(t_{n+1}-t_n)}\}}
{\sigma_{T-t_n}^2}\notag\\
&\quad=
\frac{
    e^{-(t_{n+1}-t_n)}
    \left[1-e^{-2\{T-t_n-(t_{n+1}-t_n)\}}\right]
    +e^{t_{n+1}-t_n}\{1-e^{-2(t_{n+1}-t_n)}\}
}{
    1-e^{-2(T-t_n)}
}\notag\\
&\quad=e^{t_{n+1}-t_n}.\notag
\end{align}
Finally, decompose the remaining score difference as
\begin{align}
s_{T-t_n}(Y_{t_n})
-\hat s_{T-t_n}(\hat Y_{t_n})
=
\bigl\{
    s_{T-t_n}(Y_{t_n})
    -\hat s_{T-t_n}(Y_{t_n})
\bigr\}
+
\bigl\{
    \hat s_{T-t_n}(Y_{t_n})
    -\hat s_{T-t_n}(\hat Y_{t_n})
\bigr\}.\notag
\end{align}
Substituting these identities into the preceding difference formula
gives
\begin{align}
    Y_{t_{n+1}}-\hat Y_{t_{n+1}}
    ={}&e^{t_{n+1}-t_n}\bigl(Y_{t_n}-\hat Y_{t_n}\bigr)\notag\\
    &+e^{t_{n+1}-t_n}\{1-e^{-2(t_{n+1}-t_n)}\}
      \bigl\{s_{T-t_n}(Y_{t_n})-
      \hat s_{T-t_n}(Y_{t_n})\bigr\}\notag\\
    &+e^{t_{n+1}-t_n}\{1-e^{-2(t_{n+1}-t_n)}\}
      \bigl\{\hat s_{T-t_n}(Y_{t_n})-
      \hat s_{T-t_n}(\hat Y_{t_n})\bigr\}
      +R_n.
\label{eq:error-recursion-vector}
\end{align}
Let $S_n$ denote the first three terms on the right-hand side. Since
$\bigl(\mu_{T-t}(Y_t)\bigr)_{0\le t\le t_N}$ is a martingale and the
kernel in \eqref{eq:Rn} is deterministic,
\begin{align}
\E\left[R_n\mid\mathcal F_{t_n}\right]
={}&\int_{t_n}^{t_{n+1}}
 e^{-(t_{n+1}-u)}
 \frac{\sigma_{T-t_{n+1}}^2}{\sigma_{T-u}^2}
 \frac{2e^{-(T-u)}}{\sigma_{T-u}^2}\notag\\
&\times
\E\left[
    \mu_{T-u}(Y_u)-\mu_{T-t_n}(Y_{t_n})
    \,\middle|\,
    \mathcal F_{t_n}
\right]\dd u\notag\\
={}&0.\notag
\end{align}
Thus,
\begin{equation}
    \E\left[R_n\mid\mathcal F_{t_n}\right]=0.
    \label{eq:Rn-centered}
\end{equation}
Every term defining $S_n$ is a deterministic function of
$Y_{t_n}$ and $\hat Y_{t_n}$. Hence $S_n$ is
$\mathcal F_{t_n}$-measurable. Equation~\eqref{eq:error-recursion-vector}
can be written as
\[
    Y_{t_{n+1}}-\hat Y_{t_{n+1}}=S_n+R_n.
\]
Expanding the squared norm and taking expectations gives
\begin{align}
\E\left[
    \norm{Y_{t_{n+1}}-\hat Y_{t_{n+1}}}^2
\right]
=
\E\left[\norm{S_n}^2\right]
+2\E\left[\ip{S_n}{R_n}\right]
+\E\left[\norm{R_n}^2\right].\notag
\end{align}
The role of the $\mathcal F_{t_n}$-measurability of $S_n$ is to
allow it to be pulled outside the inner conditional expectation:
\begin{align}
\E\left[\ip{S_n}{R_n}\right]
&=
\E\left[
    \E\left[
        \ip{S_n}{R_n}
        \,\middle|\,
        \mathcal F_{t_n}
    \right]
\right]\notag\\
&=
\E\left[
    \ip{
        S_n
    }{
        \E\left[R_n\mid\mathcal F_{t_n}\right]
    }
\right]\notag\\
&=0,\notag
\end{align}
where the last equality follows from \eqref{eq:Rn-centered}.
Consequently,
\begin{equation}
    \E\left[\norm{Y_{t_{n+1}}-\hat Y_{t_{n+1}}}^2\right]
    =\E\left[\norm{S_n}^2\right]+\E\left[\norm{R_n}^2\right].
    \label{eq:orthogonal-split}
\end{equation}
We next estimate the first term on the right-hand side of
\eqref{eq:orthogonal-split}. Recall from
\eqref{eq:error-recursion-vector} that
\begin{align}
S_n={}&e^{t_{n+1}-t_n}\bigl(Y_{t_n}-\hat Y_{t_n}\bigr)\notag\\
&+e^{t_{n+1}-t_n}\{1-e^{-2(t_{n+1}-t_n)}\}
  \bigl\{s_{T-t_n}(Y_{t_n})
  -\hat s_{T-t_n}(Y_{t_n})\bigr\}\notag\\
&+e^{t_{n+1}-t_n}\{1-e^{-2(t_{n+1}-t_n)}\}
  \bigl\{\hat s_{T-t_n}(Y_{t_n})
  -\hat s_{T-t_n}(\hat Y_{t_n})\bigr\}.\notag
\end{align}
Because $t_n\le t_N=T-\delta$, we have
$T-t_n\in[\delta,T]$. Hence Assumption~\ref{ass:lipschitz},
applied at time $r=T-t_n$ and at the two spatial points
$Y_{t_n}$ and $\hat Y_{t_n}$, yields
\[
\norm{
    \hat s_{T-t_n}(Y_{t_n})
    -\hat s_{T-t_n}(\hat Y_{t_n})
}
\le
L(T-t_n)\norm{Y_{t_n}-\hat Y_{t_n}}.
\]
Using this bound and the triangle inequality in the preceding
expression for $S_n$, we obtain the following pathwise estimate:
\begin{align}
\norm{S_n}
\le{}&
e^{t_{n+1}-t_n}
\left[1+L(T-t_n)\{1-e^{-2(t_{n+1}-t_n)}\}\right]
\norm{Y_{t_n}-\hat Y_{t_n}}\notag\\
&+e^{t_{n+1}-t_n}\{1-e^{-2(t_{n+1}-t_n)}\}
\norm{
    s_{T-t_n}(Y_{t_n})
    -\hat s_{T-t_n}(Y_{t_n})
}.\notag
\end{align}
We now take the $L^2$ norm of both sides. Minkowski's inequality
gives
\begin{align}
\left(\E\left[\norm{S_n}^2\right]\right)^{1/2}
\le{}&
e^{t_{n+1}-t_n}
\left[1+L(T-t_n)\{1-e^{-2(t_{n+1}-t_n)}\}\right]
\left(
    \E\left[
        \norm{Y_{t_n}-\hat Y_{t_n}}^2
    \right]
\right)^{1/2}\notag\\
&+e^{t_{n+1}-t_n}\{1-e^{-2(t_{n+1}-t_n)}\}
\left(
    \E\left[
        \norm{
            s_{T-t_n}(Y_{t_n})
            -\hat s_{T-t_n}(Y_{t_n})
        }^2
    \right]
\right)^{1/2}.\notag
\end{align}
Finally, since the exact reverse process was defined by
$Y_t=X_{T-t}$, we have
$Y_{t_n}\overset{\mathrm d}=X_{T-t_n}$. Therefore
\begin{align}
&\E\left[
    \norm{
        s_{T-t_n}(Y_{t_n})
        -\hat s_{T-t_n}(Y_{t_n})
    }^2
\right]\notag\\
&\qquad=
\E\left[
    \norm{
        s_{T-t_n}(X_{T-t_n})
        -\hat s_{T-t_n}(X_{T-t_n})
    }^2
\right]
\le\varepsilon_n^2,\notag
\end{align}
where the last inequality is precisely
Assumption~\ref{ass:score-error} for the index $n$. Substituting this
estimate into the preceding Minkowski bound proves
\begin{align}
    \left(\E\left[\norm{S_n}^2\right]\right)^{1/2}
    \le{}&e^{t_{n+1}-t_n}
    \left[1+L(T-t_n)\{1-e^{-2(t_{n+1}-t_n)}\}\right]
    \left(\E\left[\norm{Y_{t_n}-\hat Y_{t_n}}^2\right]\right)^{1/2}\notag\\
    &+e^{t_{n+1}-t_n}\{1-e^{-2(t_{n+1}-t_n)}\}\varepsilon_n.
\label{eq:Sn-bound}
\end{align}

To propagate the one-step estimate over the whole time grid, we use
the following discrete quadratic recursion.
\begin{lemma}[quadratic discrete Gronwall inequality]
\label{lem:discrete}
Let $x_n,B_n,C_n\ge0$ and $A_n\in\R$ satisfy
\begin{equation}
    x_{n+1}^2\le(e^{A_n}x_n+B_n)^2+C_n^2.
    \label{eq:discrete-hyp}
\end{equation}
Then
\begin{align}
    x_{n+1}\le{}&e^{\sum_{i=0}^nA_i}x_0
    +\sum_{j=0}^n e^{\sum_{i=j+1}^nA_i}B_j\notag\\
    &+\left\{\sum_{j=0}^n e^{2\sum_{i=j+1}^nA_i}C_j^2\right\}^{1/2}.
\label{eq:discrete-conclusion}
\end{align}
\end{lemma}

The proof is deferred to \Cref{app:discrete}.

Applying Lemma~\ref{lem:discrete} to \eqref{eq:orthogonal-split}--\eqref{eq:Sn-bound} and using that each stability factor is at least one gives
\begin{align}
    \left(\E\left[\norm{Y_{t_N}-\hat Y_{t_N}}^2\right]\right)^{1/2}
    \le \mathcal G_N\Biggl[
       &\left(\E\left[\norm{Y_0-\hat Y_0}^2\right]\right)^{1/2}
       +\sum_{j=0}^{N-1}e^{t_{j+1}-t_j}
       \{1-e^{-2(t_{j+1}-t_j)}\}\varepsilon_j\notag\\
       &+\left\{\sum_{j=0}^{N-1}\E\left[\norm{R_j}^2\right]\right\}^{1/2}
    \Biggr].
\label{eq:global-recursion}
\end{align}

\subsection{Initialization and discretization errors}

The error decomposition is now complete. The first term on the
right-hand side of \eqref{eq:global-recursion} is the remaining
initialization error, which we estimate as follows. Let
$Z\sim\N(0,\Id)$ be independent of $X_0$. Then
\[
    e^{-T}X_0+\sigma_TZ\sim q_T,
    \qquad
    Z\sim\gamma_d.
\]
Thus,
$\bigl(e^{-T}X_0+\sigma_TZ,Z\bigr)$ is an admissible
coupling of $q_T$ and $\gamma_d$. On the other hand,
\eqref{eq:optimal-initial-cost} states that the actual initial
coupling is optimal and therefore gives the first equality below.
Combining these two facts yields
\begin{align}
    \E\left[\norm{Y_0-\hat Y_0}^2\right]
    &=
    W_2^2(q_T,\gamma_d)\notag\\
    &\le
    \E\left[
        \norm{e^{-T}X_0+\sigma_TZ-Z}^2
    \right]\notag\\
    &=
    e^{-2T}\E\left[\norm{X_0}^2\right]
    +(1-\sigma_T)^2d.
\label{eq:delta0-bound}
\end{align}
The explicit Gaussian coupling above is used only to bound the
optimal transportation cost. The actual initial coupling employed
in the synchronous construction is the optimal-map coupling
\eqref{eq:explicit-initial-coupling}.

Having controlled the initialization term in
\eqref{eq:global-recursion}, we now turn to the temporal
discretization term. Rather than bounding the posterior-covariance
decrement separately on each time interval, we count its total change
through a single information budget.

Define
\begin{equation}
    V(u):=\E\left[\Tr\Sigma_u(X_u)\right],
    \qquad u\in[\delta,T].
    \label{eq:V-def}
\end{equation}
By \eqref{eq:martingale-isometry} and
$Y_t\overset{\mathrm d}=X_{T-t}$, the function $V$ is nondecreasing.
The martingale isometry gives the following two equivalent
representations of the temporal discretization term:
\begin{align}
    &\sum_{n=0}^{N-1}
    \int_{t_n}^{t_{n+1}}
    \frac{1-\sigma_{T-t}^2}{\sigma_{T-t}^4}
    \E\left[
       \norm{
          \mu_{T-t}(Y_t)
          -\mu_{T-t_n}(Y_{t_n})
       }^2
    \right]\dd t\notag\\
    &=
    \sum_{n=0}^{N-1}
    \int_{t_n}^{t_{n+1}}
    \frac{1-\sigma_{T-t}^2}{\sigma_{T-t}^4}
    \bigl\{V(T-t_n)-V(T-t)\bigr\}\dd t.
\label{eq:T2-definition}
\end{align}
The second equality is exactly the martingale isometry
\eqref{eq:martingale-isometry}; thus the quantity being accumulated is
the posterior-covariance decrement, not a collection of unrelated
local errors.

\begin{lemma}[Information-theoretic accumulated remainder]
\label{lem:remainder}
Under the assumptions of \Cref{thm:main}, the common value in
\eqref{eq:T2-definition} has the exact representation
\begin{equation}
    \frac12
    \sum_{n=0}^{N-1}
    \int_{(T-t_{n+1},T-t_n]}
    \bigl\{\gamma(T-t_{n+1})-\gamma(v)\bigr\}\dd V(v).
    \label{eq:T2-Stieltjes}
\end{equation}
Moreover, this expression satisfies
\begin{align}
    &\frac12
    \sum_{n=0}^{N-1}
    \int_{(T-t_{n+1},T-t_n]}
    \bigl\{\gamma(T-t_{n+1})-\gamma(v)\bigr\}\dd V(v)\notag\\
    &\quad\le
    \frac12\left(\frac{e^{2\kappa}}{1-\kappa}-1\right)
    \int_\delta^T\gamma(v)\dd V(v)
    \label{eq:T2-SNR-budget}\\
    &\le
    \left(\frac{e^{2\kappa}}{1-\kappa}-1\right)
    I(X_0;X_\delta)
    \label{eq:T2-mutual-information}\\
    &\le
    \left(\frac{e^{2\kappa}}{1-\kappa}-1\right)
    \inf_{r>0}
    \left\{
       \log N_{\mathrm{cov}}(\mathcal X,\norm{\cdot},r)
       +\frac{\gamma(\delta)r^2}{2}
    \right\}.
    \label{eq:T2-entropy-envelope}
\end{align}
Consequently, the local remainders in \eqref{eq:Rn} satisfy
\begin{equation}
    \sum_{n=0}^{N-1}\E\left[\norm{R_n}^2\right]
    \le C\kappa k\log k.
    \label{eq:remainder-bound}
\end{equation}
\end{lemma}

The proof is deferred to \Cref{app:remainder}. Its four steps are
visible in \eqref{eq:T2-Stieltjes}--\eqref{eq:T2-entropy-envelope}.
First, the identity
\[
    \frac{1-\sigma_u^2}{\sigma_u^4}
    =-\frac12\gamma'(u)
\]
converts \eqref{eq:T2-definition} into an exact
Lebesgue--Stieltjes integral. Second, the relative mesh condition
implies
\[
    \gamma(T-t_{n+1})-\gamma(v)
    \le
    \left(\frac{e^{2\kappa}}{1-\kappa}-1\right)\gamma(v),
    \qquad v\in[T-t_{n+1},T-t_n].
\]
Third, the I--MMSE identity bounds the resulting SNR-weighted total
variation by $2I(X_0;X_\delta)$. Finally, a covering argument bounds
this mutual information by the metric-entropy envelope in
\eqref{eq:T2-entropy-envelope}. Taking
$r=\varepsilon_0=k^{-C_0}$, using
$2C_0>C_\delta$, and observing that
\[
    \frac{e^{2\kappa}}{1-\kappa}-1
    \lesssim\kappa,
    \qquad \kappa\le0.9,
\]
yields
\eqref{eq:remainder-bound}. This argument introduces neither a
separate factor $T$ nor a term of order $\kappa^2N$.

\subsection{Completion of the proof}

Combining \eqref{eq:w2-coupling}, \eqref{eq:global-recursion},
\eqref{eq:delta0-bound}, and Lemma~\ref{lem:remainder} gives the three
error terms in \eqref{eq:post-early-stopping-bound}. Finally,
$1+x\le e^x$, $1-e^{-2(t_{i+1}-t_i)}
\le2(t_{i+1}-t_i)$, and
$\sum_{i=0}^{N-1}(t_{i+1}-t_i)=t_N$ imply
\[
\mathcal G_N
\le\exp\left\{t_N
+2\sum_{i=0}^{N-1}(t_{i+1}-t_i)L(T-t_i)\right\},
\]
which is the exponential stability factor appearing in
\eqref{eq:post-early-stopping-bound}. This completes the proof of
\Cref{thm:main}.

\section{Discussion}
\label{sec:discussion}

\paragraph{Why the posterior mean is the right state variable.}
A direct discretization estimate for the score can introduce derivatives of the score and ambient-dimensional moment terms. Tweedie's formula replaces the nonlinear score contribution by the posterior mean. Its reverse-time martingale structure turns temporal increments into differences of posterior variances. The identity $(1-\sigma_u^2)/\sigma_u^4=-\gamma'(u)/2$ then makes the DDPM parameterization especially compatible with an SNR-clock analysis: the entire covariance decrement becomes one Stieltjes integral rather than a sum of separately bounded local variations.

\paragraph{Why mutual information removes the extra logarithm.}
The adaptive mesh pays on each interval only the relative SNR change
\[
    \frac{e^{2\kappa}}{1-\kappa}-1=O(\kappa).
\]
The I--MMSE identity then
converts the SNR-weighted total variation of the MMSE curve into the
single information budget $I(X_0;X_\delta)$. Metric entropy controls
this budget directly, so the same posterior-covariance change is never
charged once through the elapsed time $T$ and again through the number
of intervals $N$. This is why both the $T$ factor and the
$\kappa^2N$ term disappear, and why the KL iteration complexity loses
one logarithmic factor.

\paragraph{Role of the exponential integrator.}
The interpolation \eqref{eq:interpolation} exactly reproduces the DDPM coefficients at grid points. This exact correspondence is useful for interpreting the theorem as a guarantee for the original DDPM sampler rather than for a nearby Euler--Maruyama scheme. A standard Euler discretization can also be analyzed, but it produces additional local errors from the singular linear coefficient near the data endpoint and no longer matches \eqref{eq:ddpm-update} exactly.

\paragraph{Limitations.}
The Wasserstein theorem requires bounded support, a one-scale
metric-entropy condition, a polynomial relation between $d$ and $k$,
and a Lipschitz learned score. The stability factor may be large when
the learned score has a large Lipschitz modulus. The one-scale
information estimate also requires compatibility between the covering
resolution and early stopping, expressed here by
$2C_0>C_\delta$. More generally, it is enough that
$\gamma(\delta)\varepsilon_0^2\lesssim k\log k$; for substantially
smaller $\delta$, one would need covering information at finer scales.
In addition, ambient dimension remains in the initialization and
early-stopping terms. Obtaining a completely intrinsic-dimensional
$W_2$ bound under minimal score regularity remains open.

\paragraph{Possible extensions.}
Natural directions include replacing global Lipschitz continuity by one-sided or average regularity, extending the posterior-covariance argument to unbounded distributions with concentration assumptions, optimizing the nonuniform time mesh, and analyzing the robustness of the intrinsic-dimensional bound under perturbations of the DDPM coefficients.

\clearpage
\appendix

\begin{center}
    {\LARGE\bfseries Appendix}
\end{center}
\vspace{1.5em}

\section{Proof of the DDPM interpolation result}
\label{app:interpolation-proof}

\begin{proof}[Proof of Proposition~\ref{prop:interpolation}]
Applying the variation-of-constants formula on $[t_n,t_{n+1}]$ gives
\begin{align}
\hat Y_{t_{n+1}}
={}&e^{-(t_{n+1}-t_n)}
\frac{\sigma_{T-t_{n+1}}^2}{\sigma_{T-t_n}^2}\hat Y_{t_n}
+\frac{e^{-(T-t_{n+1})}\{1-e^{-2(t_{n+1}-t_n)}\}}{\sigma_{T-t_n}^2}
 \hat\mu_{T-t_n}(\hat Y_{t_n})\notag\\
&+\sqrt{2}\int_{t_n}^{t_{n+1}}
 e^{-(t_{n+1}-u)}
 \frac{\sigma_{T-t_{n+1}}^2}{\sigma_{T-u}^2}\,\dd B_u.
\notag
\end{align}
Substituting \eqref{eq:muhat} and using
\[
 e^{-(t_{n+1}-t_n)}
 \frac{\sigma_{T-t_{n+1}}^2}{\sigma_{T-t_n}^2}
 +\frac{e^{t_{n+1}-t_n}\{1-e^{-2(t_{n+1}-t_n)}\}}
 {\sigma_{T-t_n}^2}
 =e^{t_{n+1}-t_n}
\]
shows that the first two deterministic terms in this representation
equal
\begin{align}
    \frac{1}{\sqrt{\alpha_n}}
    \left\{
        \hat Y_{t_n}
        +(1-\alpha_n)
        \hat s_{T-t_n}(\hat Y_{t_n})
    \right\}.
    \notag
\end{align}
Let $G_n$ denote the stochastic integral in this representation,
including its factor $\sqrt{2}$.
A direct It\^o-isometry calculation gives
\begin{align}
    \E\left[G_nG_n^\top\right]
    &=
    \frac{(1-\alpha_n)(1-\bar\alpha_{n+1})}
         {1-\bar\alpha_n}\,\Id.
    \notag
\end{align}
Therefore, with
\begin{align}
    Z_n
    &:=
    \sqrt{
        \frac{1-\bar\alpha_n}
             {(1-\alpha_n)(1-\bar\alpha_{n+1})}
    }\,G_n,
    \notag
\end{align}
each $Z_n$ is standard Gaussian. Since the $G_n$ are stochastic
integrals over disjoint Brownian time intervals, the vectors $Z_n$
are independent of one another and of $\hat Y_{t_0}$. Hence the
grid-point values of $\hat Y$ obey exactly the same recursion as
\eqref{eq:ddpm-update}. Moreover,
$\hat Y_{t_0}$ and $\widetilde Y_0$ have the same distribution, and
each is independent of its corresponding Gaussian innovations.
It follows inductively that
\[
    (\widetilde Y_0,\ldots,\widetilde Y_N)
    \overset{\mathrm d}{=}
    (\hat Y_{t_0},\ldots,\hat Y_{t_N}),
\]
which proves the proposition.
\end{proof}

\section{Proof of the reverse posterior martingale}
\label{app:martingale}

\begin{proof}[Proof of Lemma~\ref{lem:martingale}]
By the construction in \Cref{sec:proof},
\[
    \mathcal F_t
    =
    \mathcal G_t
    =
    \sigma\{Y_s:0\le s\le t\}.
\]
Since $Y_s=X_{T-s}$ almost surely,
\[
    \mathcal F_t
    =
    \sigma\{X_u:T-t\le u\le T\}.
\]
The Markov property of the forward Ornstein--Uhlenbeck process implies
\[
    \E\left[X_0\mid\mathcal F_t\right]
    =
    \E\left[X_0\mid X_{T-t}\right]
    =
    \mu_{T-t}(Y_t).
\]
Because $X_0$ is square-integrable, this identity proves that
$\bigl(\mu_{T-t}(Y_t)\bigr)_{0\le t\le t_N}$ is a square-integrable
martingale.

For $0\le s\le t\le t_N$, the orthogonality of conditional
expectations gives
\begin{align}
    \E\left[\norm{X_0-\mu_{T-s}(Y_s)}^2\right]
    ={}&
    \E\left[\norm{X_0-\mu_{T-t}(Y_t)}^2\right]\notag\\
    &+
    \E\left[
        \norm{\mu_{T-t}(Y_t)-\mu_{T-s}(Y_s)}^2
    \right].\notag
\end{align}
For every $r\in[0,t_N]$, the definition of conditional covariance
yields
\[
    \E\left[\norm{X_0-\mu_{T-r}(Y_r)}^2\right]
    =
    \E\left[\Tr\Sigma_{T-r}(Y_r)\right].
\]
Substitution with $r=s$ and $r=t$ proves
\eqref{eq:martingale-isometry}.
\end{proof}

\section{Proof of the discrete quadratic recursion}
\label{app:discrete}

\begin{proof}[Proof of Lemma~\ref{lem:discrete}]
Define two nonnegative sequences $(p_n)_{n\ge0}$ and
$(q_n)_{n\ge0}$ by
\[
    p_0:=x_0,
    \qquad
    q_0:=0,
\]
and, for $n\ge0$,
\begin{equation}
    p_{n+1}:=e^{A_n}p_n+B_n,
    \qquad
    q_{n+1}:=\left(e^{2A_n}q_n^2+C_n^2\right)^{1/2}.
    \label{eq:pq-recursion}
\end{equation}
We first show by induction that
\begin{equation}
    x_n\le p_n+q_n
    \qquad\text{for every }n\ge0.
    \label{eq:pq-majorization}
\end{equation}
The claim is immediate for $n=0$. Suppose that it holds for some
$n$. Since all terms are nonnegative, \eqref{eq:discrete-hyp} and
the induction hypothesis imply
\begin{align}
    x_{n+1}
    &\le
    \left\{
        \bigl(e^{A_n}x_n+B_n\bigr)^2+C_n^2
    \right\}^{1/2}\notag\\
    &\le
    \left\{
        \bigl(e^{A_n}(p_n+q_n)+B_n\bigr)^2+C_n^2
    \right\}^{1/2}.\notag
\end{align}
Write the last expression as the Euclidean norm of the sum
\[
    \bigl(e^{A_n}p_n+B_n,0\bigr)
    +
    \bigl(e^{A_n}q_n,C_n\bigr).
\]
Minkowski's inequality in $\R^2$ therefore gives
\begin{align}
    x_{n+1}
    &\le
    e^{A_n}p_n+B_n
    +\left(e^{2A_n}q_n^2+C_n^2\right)^{1/2}\notag\\
    &=p_{n+1}+q_{n+1}.\notag
\end{align}
This proves \eqref{eq:pq-majorization}.

It remains to solve the two recursions in
\eqref{eq:pq-recursion}. Iterating the first one yields
\begin{equation}
    p_{n+1}
    =
    e^{\sum_{i=0}^nA_i}x_0
    +\sum_{j=0}^n
        e^{\sum_{i=j+1}^nA_i}B_j.
    \label{eq:p-explicit}
\end{equation}
Similarly, squaring the second recursion gives
$q_{n+1}^2=e^{2A_n}q_n^2+C_n^2$. Since $q_0=0$, iteration gives
\begin{equation}
    q_{n+1}^2
    =
    \sum_{j=0}^n
        e^{2\sum_{i=j+1}^nA_i}C_j^2.
    \label{eq:q-explicit}
\end{equation}
Combining \eqref{eq:pq-majorization},
\eqref{eq:p-explicit}, and \eqref{eq:q-explicit} proves
\eqref{eq:discrete-conclusion}.
\end{proof}

\section{Information-theoretic proof of the accumulated remainder bound}
\label{app:remainder}

\begin{proof}[Proof of Lemma~\ref{lem:remainder}]
\medskip\noindent\emph{Reduction to the posterior-covariance decrement.}
By Cauchy--Schwarz, the mesh condition, and the deterministic kernel
bound in \eqref{eq:Rn},
\begin{equation}
    \E\left[\norm{R_n}^2\right]
    \le 4\int_{t_n}^{t_{n+1}}
       \frac{1-\sigma_{T-t}^2}{\sigma_{T-t}^4}
       \E\left[
          \norm{
             \mu_{T-t}(Y_t)-\mu_{T-t_n}(Y_{t_n})
          }^2
       \right]\dd t.
    \label{eq:Rn-pre}
\end{equation}
Summing \eqref{eq:Rn-pre} over $n$ and using
\eqref{eq:T2-definition} gives
\begin{align}
    \sum_{n=0}^{N-1}\E\left[\norm{R_n}^2\right]
    &\le
    4\sum_{n=0}^{N-1}
    \int_{t_n}^{t_{n+1}}
    \frac{1-\sigma_{T-t}^2}{\sigma_{T-t}^4}\notag\\
    &\quad\times
    \E\left[
       \norm{
          \mu_{T-t}(Y_t)-\mu_{T-t_n}(Y_{t_n})
       }^2
    \right]\dd t.
    \label{eq:Rn-by-T2}
\end{align}

\medskip\noindent\emph{Exact Stieltjes representation.}
From \eqref{eq:gamma},
\begin{equation}
    \gamma'(u)
    =-\frac{2e^{2u}}{(e^{2u}-1)^2}
    =-2\gamma(u)\{1+\gamma(u)\},
    \label{eq:gamma-derivative}
\end{equation}
and therefore
\begin{equation}
    \frac{1-\sigma_u^2}{\sigma_u^4}
    =\frac{e^{-2u}}{(1-e^{-2u})^2}
    =\gamma(u)\{1+\gamma(u)\}
    =-\frac12\gamma'(u).
    \label{eq:weight-gamma-derivative}
\end{equation}
Changing variables from reverse time $t$ to forward time $u=T-t$
in \eqref{eq:T2-definition} yields
\begin{align}
    &\sum_{n=0}^{N-1}
    \int_{t_n}^{t_{n+1}}
    \frac{1-\sigma_{T-t}^2}{\sigma_{T-t}^4}
    \bigl\{V(T-t_n)-V(T-t)\bigr\}\dd t\notag\\
    &\quad=
    \sum_{n=0}^{N-1}
    \int_{T-t_{n+1}}^{T-t_n}
    \left\{-\frac12\gamma'(u)\right\}
    \{V(T-t_n)-V(u)\}\dd u.
    \label{eq:T2-forward-time}
\end{align}

The function $V$ is nondecreasing and satisfies
$0\le V(u)\le\E\left[\norm{X_0}^2\right]$. It therefore induces a
finite Lebesgue--Stieltjes measure, denoted by $\dd V(v)$, through
\[
    \int_{(a,b]}\dd V(v)=V(b)-V(a).
\]
For $u\in[T-t_{n+1},T-t_n]$,
\[
    V(T-t_n)-V(u)=\int_{(u,T-t_n]}\dd V(v).
\]
All integrands are nonnegative. Tonelli's theorem thus permits the
order of integration to be exchanged:
\begin{align}
&\int_{T-t_{n+1}}^{T-t_n}
 \left\{-\frac12\gamma'(u)\right\}
 \left\{\int_{(u,T-t_n]}\dd V(v)\right\}\dd u
 \notag\\
&\quad=
\int_{(T-t_{n+1},T-t_n]}
\left\{
   \int_{T-t_{n+1}}^v
   \left\{-\frac12\gamma'(u)\right\}\dd u
\right\}\dd V(v)
\notag\\
&\quad=
\frac12
\int_{(T-t_{n+1},T-t_n]}
\{\gamma(T-t_{n+1})-\gamma(v)\}\dd V(v).
\label{eq:Tonelli-Stieltjes}
\end{align}
Substitution into \eqref{eq:T2-forward-time}, followed by summation
over $n$, proves the exact identity \eqref{eq:T2-Stieltjes}. Each
increment $\dd V(v)$ is counted exactly once.

\medskip\noindent\emph{Relative SNR control on one time interval.}
For every $r>0$,
\begin{equation}
    -\frac{\dd}{\dd r}\log\gamma(r)
    =\frac{2}{1-e^{-2r}}
    \le2+\frac1r.
    \label{eq:log-SNR-derivative}
\end{equation}
Indeed,
\[
    \frac{1}{1-e^{-2r}}
    =1+\frac{1}{e^{2r}-1}
    \le1+\frac{1}{2r},
\]
where the last inequality uses $e^{2r}-1\ge2r$.

Fix $n$ and $v\in[T-t_{n+1},T-t_n]$, and set
$a=T-t_{n+1}$ and $b=T-t_n$. By \eqref{eq:mesh-assumption},
\[
    b-a=t_{n+1}-t_n\le\kappa\min\{1,b\}.
\]
If $b\le1$, this gives $a\ge(1-\kappa)b$. If $b>1$, then
$a\ge b-\kappa\ge(1-\kappa)b$. Consequently,
\begin{equation}
    a\ge(1-\kappa)b\ge(1-\kappa)v,
    \qquad
    v-a\le\kappa.
    \label{eq:relative-step-geometry}
\end{equation}
Integrating \eqref{eq:log-SNR-derivative} and using
\eqref{eq:relative-step-geometry}, we obtain
\begin{align}
    \log\frac{\gamma(a)}{\gamma(v)}
    &=
    \int_a^v\frac{2}{1-e^{-2r}}\dd r
    \notag\\
    &\le
    2(v-a)+\log\frac{v}{a}
    \notag\\
    &\le
    2\kappa+\log\frac{1}{1-\kappa}.
    \label{eq:SNR-log-ratio}
\end{align}
After exponentiating,
\[
    \frac{\gamma(a)}{\gamma(v)}
    \le\frac{e^{2\kappa}}{1-\kappa}.
\]
Hence
\begin{equation}
    \gamma(T-t_{n+1})-\gamma(v)
    \le
    \left(\frac{e^{2\kappa}}{1-\kappa}-1\right)\gamma(v),
    \qquad v\in[T-t_{n+1},T-t_n].
    \label{eq:relative-SNR-control}
\end{equation}
Applying \eqref{eq:relative-SNR-control} in
\eqref{eq:T2-Stieltjes}, and observing that the intervals
$(T-t_{n+1},T-t_n]$ partition $(\delta,T]$, proves
\eqref{eq:T2-SNR-budget}. Moreover, for every fixed
$\kappa_0<1$, the mean-value theorem gives
\begin{equation}
    \frac{e^{2\kappa}}{1-\kappa}-1
    \le
    \kappa e^{2\kappa_0}
    \left\{
       \frac{2}{1-\kappa_0}
       +\frac{1}{(1-\kappa_0)^2}
    \right\},
    \qquad 0<\kappa\le\kappa_0.
    \label{eq:relative-mesh-linear}
\end{equation}
In particular,
\[
    \frac{e^{2\kappa}}{1-\kappa}-1
    \lesssim\kappa
    \qquad\text{for }\kappa\le0.9.
\]

\medskip\noindent\emph{The I--MMSE information budget.}
Let $G\sim\N(0,\Id)$ be independent of $X_0$. For $\lambda\ge0$,
define the Gaussian channel, its MMSE, and its mutual information by
\begin{align}
    Z_\lambda&:=\sqrt{\lambda}\,X_0+G,\notag\\
    m(\lambda)
    &:=
    \E\left[
       \norm{
          X_0-\E\left[X_0\mid Z_\lambda\right]
       }^2
    \right],\notag\\
    \mathcal I(\lambda)&:=I(X_0;Z_\lambda).
\label{eq:Gaussian-channel-defs}
\end{align}
Since
\[
    \frac{X_u}{\sigma_u}
    \overset{\mathrm d}=
    \sqrt{\gamma(u)}\,X_0+G,
\]
and multiplication of the observation by a nonzero constant does not
change its generated sigma-field,
\begin{equation}
    V(u)=m(\gamma(u)),
    \qquad
    I(X_0;X_u)=\mathcal I(\gamma(u)).
    \label{eq:OU-channel-identification}
\end{equation}

The function $m$ is nonincreasing. To see this, take
$0\le\lambda_1\le\lambda_2$ and generate
\[
    \widetilde Z_{\lambda_1}
    :=
    \sqrt{\frac{\lambda_1}{\lambda_2}}Z_{\lambda_2}
    +\sqrt{1-\frac{\lambda_1}{\lambda_2}}G',
\]
where $G'\sim\N(0,\Id)$ is independent. Then
$X_0\to Z_{\lambda_2}\to\widetilde Z_{\lambda_1}$ is a Markov chain
and $\widetilde Z_{\lambda_1}$ has the same conditional distribution
given $X_0$ as $Z_{\lambda_1}$. The $L^2$ projection property of
conditional expectation therefore gives
$m(\lambda_2)\le m(\lambda_1)$.

The vector I--MMSE identity \citep{guo2005mutual} states that
$\mathcal I$ is absolutely continuous and
\begin{equation}
    \mathcal I'(\lambda)=\frac12m(\lambda)
    \quad\text{for almost every }\lambda\ge0.
    \label{eq:I-MMSE}
\end{equation}
Thus, for $0\le a\le b$,
\begin{equation}
    \int_a^b m(\lambda)\dd\lambda
    =2\{\mathcal I(b)-\mathcal I(a)\}.
    \label{eq:I-MMSE-integral}
\end{equation}

Both $V$ and $\gamma$ are continuous and of bounded variation on
$[\delta,T]$. Stieltjes integration by parts, followed by the change
of variables $\lambda=\gamma(u)$, gives
\begin{align}
    \int_\delta^T\gamma(u)\dd V(u)
    ={}&
    \gamma(T)m(\gamma(T))
    -\gamma(\delta)m(\gamma(\delta))
    \notag\\
    &+
    2\{\mathcal I(\gamma(\delta))-\mathcal I(\gamma(T))\}.
    \label{eq:integrated-budget-exact}
\end{align}
Because $m$ is nonincreasing and $\mathcal I(0)=0$,
\begin{equation}
    \gamma(T)m(\gamma(T))
    \le
    \int_0^{\gamma(T)}m(\lambda)\dd\lambda
    =2\mathcal I(\gamma(T)).
    \label{eq:low-SNR-endpoint}
\end{equation}
Substituting \eqref{eq:low-SNR-endpoint} into
\eqref{eq:integrated-budget-exact} and dropping the nonnegative term
$\gamma(\delta)m(\gamma(\delta))$ yields
\begin{equation}
    \int_\delta^T\gamma(u)\dd V(u)
    \le2\mathcal I(\gamma(\delta))
    =2I(X_0;X_\delta).
    \label{eq:integrated-budget-information}
\end{equation}
Combining \eqref{eq:T2-SNR-budget} and
\eqref{eq:integrated-budget-information} proves
\eqref{eq:T2-mutual-information}. Notice that the terminal time $T$
does not appear on the right-hand side: the low-SNR information below
$\gamma(T)$ absorbs the endpoint contribution in
\eqref{eq:integrated-budget-exact}.

\medskip\noindent\emph{Metric entropy controls the mutual information.}
Fix $r>0$, let
$M=N_{\mathrm{cov}}(\mathcal X,\norm{\cdot},r)<\infty$, and choose an
$r$-net $\{x_1^*,\ldots,x_M^*\}$ of $\mathcal X$. Let $J=J(X_0)$ be
the smallest index of a nearest net point. Then $J$ is a measurable
function of $X_0$ and
\begin{equation}
    \norm{X_0-x_J^*}\le r
    \qquad\text{almost surely}.
    \label{eq:quantization-residual}
\end{equation}
For $Z_\lambda=\sqrt{\lambda}X_0+G$, the mutual-information chain rule
gives
\begin{align}
    I(X_0;Z_\lambda)
    &=
    I(J;Z_\lambda)+I(X_0;Z_\lambda\mid J)
    \notag\\
    &\le
    \log M+I(X_0;Z_\lambda\mid J).
    \label{eq:information-chain-net}
\end{align}
Conditional on $J=j$, put $R_j=X_0-x_j^*$ and
$K_j=\Cov(R_j\mid J=j)$. Translating the channel output by
$\sqrt{\lambda}x_j^*$ does not change mutual information. The
maximum-entropy property of the Gaussian distribution
\citep{cover2006elements} and
$\log\det(\Id+A)\le\Tr A$ therefore give
\begin{align}
&I(X_0;Z_\lambda\mid J=j)
\notag\\
&\quad=
I(R_j;\sqrt{\lambda}R_j+G\mid J=j)
\notag\\
&\quad\le
\frac12\log\det(\Id+\lambda K_j)
\notag\\
&\quad\le
\frac{\lambda}{2}\Tr K_j
\notag\\
&\quad\le
\frac{\lambda}{2}
\E\left[\norm{R_j}^2\mid J=j\right]
\le\frac{\lambda r^2}{2}.
\label{eq:conditional-channel-information}
\end{align}
Averaging over $j$ in \eqref{eq:information-chain-net} proves
\begin{equation}
    I(X_0;\sqrt{\lambda}X_0+G)
    \le
    \log N_{\mathrm{cov}}(\mathcal X,\norm{\cdot},r)
    +\frac{\lambda r^2}{2}.
    \label{eq:covering-information}
\end{equation}
Taking $\lambda=\gamma(\delta)$, using
\eqref{eq:OU-channel-identification}, and then taking the infimum over
$r>0$ proves \eqref{eq:T2-entropy-envelope}.

It remains to specialize this envelope to
$r=\varepsilon_0=k^{-C_0}$. Assumption~\ref{ass:entropy} gives
\[
    \log N_{\mathrm{cov}}
    (\mathcal X,\norm{\cdot},\varepsilon_0)
    \le C_{\mathrm{cov}}C_0k\log k.
\]
Also, $e^{2\delta}-1\ge2\delta$ and
$\delta\ge k^{-C_\delta}$ imply
\begin{equation}
    \frac{\gamma(\delta)\varepsilon_0^2}{2}
    \le
    \frac14 k^{C_\delta-2C_0}
    =O(1),
    \label{eq:scale-compatibility}
\end{equation}
where the last equality uses $2C_0>C_\delta$. Consequently,
$I(X_0;X_\delta)\lesssim k\log k$. Combining this result with
\eqref{eq:relative-mesh-linear},
\eqref{eq:T2-mutual-information}, and
\eqref{eq:Rn-by-T2} proves the $C\kappa k\log k$ estimate in
\eqref{eq:remainder-bound}.
\end{proof}

\section{Proof of the KL-divergence theorem}
\label{app:KL}

\begin{proof}[Proof of Theorem~\ref{thm:main-KL}]
Let $Q$ be the path law of the exact reverse process
\eqref{eq:reverse-posterior}, initialized from $q_T$, and let $P$ be
the path law of the semilinear interpolation \eqref{eq:interpolation},
initialized from $\gamma_d=\N(0,\Id)$. The endpoint marginal under
$Q$ is $q_\delta$, whereas the endpoint marginal under $P$ is
$p_{\mathrm{out}}$. The chain rule for relative entropy, Girsanov's
theorem, and data processing give
\begin{align}
    \KL(q_\delta\Vert p_{\mathrm{out}})
    \le{}&
    \KL(q_T\Vert\gamma_d)
    \notag\\
    &+
    \sum_{n=0}^{N-1}
    \int_{t_n}^{t_{n+1}}
    \frac{1-\sigma_{T-t}^2}{\sigma_{T-t}^4}
    \E\left[
       \norm{
          \mu_{T-t}(Y_t)
          -\hat\mu_{T-t_n}(Y_{t_n})
       }^2
    \right]\dd t.
    \label{eq:KL-Girsanov}
\end{align}
Here the coefficient follows directly from the common diffusion
coefficient $\sqrt2\,\Id$ and the posterior-mean drift coefficient
$2e^{-(T-t)}/\sigma_{T-t}^2$. The usual localization argument
justifies \eqref{eq:KL-Girsanov} whenever the global exponential
martingale condition is not imposed directly; see, for example,
\citet{benton2024nearly} and \citet{huang2026optimal}.

For $t\in[t_n,t_{n+1}]$, decompose the difference in
\eqref{eq:KL-Girsanov} as
\begin{align*}
&\mu_{T-t}(Y_t)-\hat\mu_{T-t_n}(Y_{t_n})\\
&\quad=
\{\mu_{T-t}(Y_t)-\mu_{T-t_n}(Y_{t_n})\}\\
&\qquad+
\{\mu_{T-t_n}(Y_{t_n})
  -\hat\mu_{T-t_n}(Y_{t_n})\}.
\end{align*}
The second term is $\mathcal F_{t_n}$-measurable, while the martingale
property in Lemma~\ref{lem:martingale} gives
\[
    \E\left[
       \mu_{T-t}(Y_t)-\mu_{T-t_n}(Y_{t_n})
       \,\middle|\,
       \mathcal F_{t_n}
    \right]=0.
\]
Therefore the cross term vanishes:
\begin{align}
&\E\left[
  \ip{
     \mu_{T-t}(Y_t)-\mu_{T-t_n}(Y_{t_n})
  }{
     \mu_{T-t_n}(Y_{t_n})
     -\hat\mu_{T-t_n}(Y_{t_n})
  }
\right]
\notag\\
&\quad=
\E\left[
  \ip{
     \E\left[
        \mu_{T-t}(Y_t)-\mu_{T-t_n}(Y_{t_n})
        \,\middle|\,
        \mathcal F_{t_n}
     \right]
  }{
     \mu_{T-t_n}(Y_{t_n})
     -\hat\mu_{T-t_n}(Y_{t_n})
  }
\right]
=0.
\label{eq:KL-cross-zero}
\end{align}
Consequently, the drift contribution in \eqref{eq:KL-Girsanov}
splits exactly into the temporal discretization term in
\eqref{eq:T2-definition} and the score-estimation term
$\mathcal T_1$, where
\begin{align}
    \mathcal T_1
    :={}&
    \sum_{n=0}^{N-1}
    \int_{t_n}^{t_{n+1}}
    \frac{1-\sigma_{T-t}^2}{\sigma_{T-t}^4}\notag\\
    &\quad\times
    \E\left[
       \norm{
          \mu_{T-t_n}(Y_{t_n})
          -\hat\mu_{T-t_n}(Y_{t_n})
       }^2
    \right]\dd t.
\label{eq:T1-definition}
\end{align}

By Tweedie's formula \eqref{eq:tweedie}, the definition
\eqref{eq:muhat}, and
$Y_{t_n}\overset{\mathrm d}=X_{T-t_n}$,
\begin{align}
&\E\left[
   \norm{
      \mu_{T-t_n}(Y_{t_n})
      -\hat\mu_{T-t_n}(Y_{t_n})
   }^2
\right]
\notag\\
&\quad=
e^{2(T-t_n)}\sigma_{T-t_n}^4
\E\left[
   \norm{
      s_{T-t_n}(X_{T-t_n})
      -\hat s_{T-t_n}(X_{T-t_n})
   }^2
\right]
\notag\\
&\quad\le
e^{2(T-t_n)}\sigma_{T-t_n}^4\varepsilon_n^2.
\label{eq:T1-score-bound}
\end{align}
Substituting \eqref{eq:T1-score-bound} into
\eqref{eq:T1-definition} proves
\begin{equation}
    \mathcal T_1\le\scoreerr^2.
    \label{eq:T1-by-score-budget}
\end{equation}
For reference, after substituting \eqref{eq:T1-score-bound} into
\eqref{eq:T1-definition}, the coefficient of $\varepsilon_n^2$ can
also be written as
\begin{equation}
    \frac{e^{2(t_{n+1}-t_n)}-1}{2}
    \frac{\sigma_{T-t_n}^2}{\sigma_{T-t_{n+1}}^2}.
    \label{eq:KL-score-weight-explicit}
\end{equation}

We next bound the initialization divergence. Since $q_T$ is the
mixture of
$\N(e^{-T}x,\sigma_T^2\Id)$ over $x\sim\pdata$, convexity of relative
entropy in its first argument gives
\begin{align}
    \KL(q_T\Vert\gamma_d)
    \le{}&
    \frac12\left[
       e^{-2T}\E\left[\norm{X_0}^2\right]
       +d\{-e^{-2T}-\log(1-e^{-2T})\}
    \right]
    \notag\\
    \le{}&
    C\left(d+\E\left[\norm{X_0}^2\right]\right)e^{-2T},
    \label{eq:initial-KL-bound}
\end{align}
where the second inequality uses $T>1$ and
$-a-\log(1-a)\le Ca^2$ for $0\le a\le e^{-2}$.

Combining \eqref{eq:KL-Girsanov}, \eqref{eq:KL-cross-zero},
\eqref{eq:T1-by-score-budget}, \eqref{eq:initial-KL-bound}, and
\eqref{eq:T2-mutual-information} gives
\begin{align}
    \KL(q_\delta\Vert p_{\mathrm{out}})
    \le{}&
    C\left(d+\E\left[\norm{X_0}^2\right]\right)e^{-2T}\notag\\
    &+\scoreerr^2
    +\left(\frac{e^{2\kappa}}{1-\kappa}-1\right)
    I(X_0;X_\delta).
\label{eq:KL-information-combined}
\end{align}
Applying the covering-information bound
\eqref{eq:covering-information} to
\eqref{eq:KL-information-combined}, taking
$r=\varepsilon_0=k^{-C_0}$, and using
\eqref{eq:scale-compatibility} and
\eqref{eq:relative-mesh-linear} proves \eqref{eq:main-KL}.

For the standard two-stage mesh,
$\kappa\asymp\{T+\log(1/\delta)\}/N$. Hence
\eqref{eq:KL-iteration-general} makes the discretization term in
\eqref{eq:main-KL} of order $\varepsilon^2$. The stated choice of
$T$ makes \eqref{eq:initial-KL-bound} of the same order, and the
assumption $\scoreerr^2\lesssim\varepsilon^2$ handles the remaining
term. Finally, Assumption~\ref{ass:support} and
\eqref{eq:poly-regime} imply
\[
    d+\E\left[\norm{X_0}^2\right]\le k^C,
    \qquad
    \log(1/\delta)\le C_\delta\log k.
\]
Thus one may take $T=O(\log(k/\varepsilon))$, which reduces
\eqref{eq:KL-iteration-general} to
\eqref{eq:KL-iteration-sharp}; \eqref{eq:KL-iteration} follows from
$\log k\le\log(k/\varepsilon)$.
\end{proof}

\section{Proof of the overall Wasserstein corollary}
\label{app:overall-W2}

\begin{proof}[Proof of Corollary~\ref{cor:overall-W2}]
The triangle inequality gives
\begin{align}
    \Wtwo\bigl(\pdata,\Law(\widetilde Y_N)\bigr)
    \le
    \Wtwo\bigl(\pdata,\Law(X_\delta)\bigr)
    +\Wtwo\bigl(\Law(X_\delta),\Law(\widetilde Y_N)\bigr).
    \label{eq:overall-error-decomposition}
\end{align}
Using the forward representation \eqref{eq:forward-marginal} with the
same $X_0$ gives
\begin{align}
    \Wtwo^2\bigl(\pdata,\Law(X_\delta)\bigr)
    &\le
    \E\left[
        \norm{X_0-e^{-\delta}X_0-\sigma_\delta Z}^2
    \right]\notag\\
    &=
    (1-e^{-\delta})^2\E\left[\norm{X_0}^2\right]
    +\sigma_\delta^2d.
    \label{eq:early-stopping}
\end{align}
Combining this estimate with \Cref{thm:main} proves
\eqref{eq:overall-W2-bound}.
\end{proof}

\begin{thebibliography}{99}

\bibitem[Anderson(1982)]{anderson1982reverse}
B. D. O. Anderson.
\newblock Reverse-time diffusion equation models.
\newblock \emph{Stochastic Processes and their Applications}, 12(3):313--326, 1982.

\bibitem[Arsenyan et al.(2025)Arsenyan, Vardanyan, and Dalalyan]{arsenyan2025assessing}
V. Arsenyan, E. Vardanyan, and A. S. Dalalyan.
\newblock Assessing the quality of denoising diffusion models in Wasserstein distance: noisy score and optimal bounds.
\newblock In \emph{Advances in Neural Information Processing Systems 38}, 2025.

\bibitem[Azangulov et al.(2024)Azangulov, Deligiannidis, and Rousseau]{azangulov2024manifold}
I. Azangulov, G. Deligiannidis, and J. Rousseau.
\newblock Convergence of diffusion models under the manifold hypothesis in high dimensions.
\newblock arXiv:2409.18804, 2024.

\bibitem[Benton et al.(2024)Benton, De Bortoli, Doucet, and Deligiannidis]{benton2024nearly}
J. Benton, V. De Bortoli, A. Doucet, and G. Deligiannidis.
\newblock Nearly $d$-linear convergence bounds for diffusion models via stochastic localization.
\newblock In \emph{International Conference on Learning Representations}, 2024.

\bibitem[Beyler and Bach(2025)]{beyler2025convergence}
E. Beyler and F. Bach.
\newblock Convergence of deterministic and stochastic diffusion-model samplers: a simple analysis in Wasserstein distance.
\newblock arXiv:2508.03210, 2025.

\bibitem[Bruno and Sabanis(2025)]{bruno2025wasserstein}
S. Bruno and S. Sabanis.
\newblock Wasserstein convergence of score-based generative models under semiconvexity and discontinuous gradients.
\newblock \emph{Transactions on Machine Learning Research}, 2025.

\bibitem[Chen et al.(2023a)Chen, Chewi, Li, Li, Salim, and Zhang]{chen2023sampling}
S. Chen, S. Chewi, J. Li, Y. Li, A. Salim, and A. R. Zhang.
\newblock Sampling is as easy as learning the score: theory for diffusion models with minimal data assumptions.
\newblock In \emph{International Conference on Learning Representations}, 2023a.

\bibitem[Chen et al.(2023b)Chen, Lee, and Lu]{chen2023improved}
H. Chen, H. Lee, and J. Lu.
\newblock Improved analysis of score-based generative modeling: user-friendly bounds under minimal smoothness assumptions.
\newblock In \emph{Proceedings of the 40th International Conference on Machine Learning}, 2023b.

\bibitem[Cover and Thomas(2006)]{cover2006elements}
T. M. Cover and J. A. Thomas.
\newblock \emph{Elements of Information Theory}.
\newblock Wiley-Interscience, second edition, 2006.

\bibitem[Efron(2011)]{efron2011tweedie}
B. Efron.
\newblock Tweedie's formula and selection bias.
\newblock \emph{Journal of the American Statistical Association}, 106(496):1602--1614, 2011.

\bibitem[Gao et al.(2025)Gao, Nguyen, and Zhu]{gao2025wasserstein}
X. Gao, H. M. Nguyen, and L. Zhu.
\newblock Wasserstein convergence guarantees for a general class of score-based generative models.
\newblock \emph{Journal of Machine Learning Research}, 26(43):1--54, 2025.

\bibitem[Gentiloni Silveri and Ocello(2025)]{gentiloni2025beyond}
M. Gentiloni Silveri and A. Ocello.
\newblock Beyond log-concavity and score regularity: improved convergence bounds for score-based generative models in $W_2$-distance.
\newblock In \emph{Proceedings of the 42nd International Conference on Machine Learning}, volume 267 of \emph{Proceedings of Machine Learning Research}, pages 55631--55656, 2025.

\bibitem[Guo et al.(2005)Guo, Shamai, and Verd\'u]{guo2005mutual}
D. Guo, S. Shamai, and S. Verd\'u.
\newblock Mutual information and minimum mean-square error in Gaussian channels.
\newblock \emph{IEEE Transactions on Information Theory}, 51(4):1261--1282, 2005.

\bibitem[Heusel et al.(2017)Heusel, Ramsauer, Unterthiner, Nessler, and Hochreiter]{heusel2017gans}
M. Heusel, H. Ramsauer, T. Unterthiner, B. Nessler, and S. Hochreiter.
\newblock GANs trained by a two time-scale update rule converge to a local Nash equilibrium.
\newblock In \emph{Advances in Neural Information Processing Systems 30}, 2017.

\bibitem[Haussmann and Pardoux(1986)]{haussmann1986time}
U. G. Haussmann and E. Pardoux.
\newblock Time reversal of diffusions.
\newblock \emph{The Annals of Probability}, 14(4):1188--1205, 1986.

\bibitem[Ho et al.(2020)Ho, Jain, and Abbeel]{ho2020ddpm}
J. Ho, A. Jain, and P. Abbeel.
\newblock Denoising diffusion probabilistic models.
\newblock In \emph{Advances in Neural Information Processing Systems 33}, pages 6840--6851, 2020.

\bibitem[Huang et al.(2026)Huang, Wei, and Chen]{huang2026optimal}
Z. Huang, Y. Wei, and Y. Chen.
\newblock Denoising diffusion probabilistic models are optimally adaptive to unknown low dimensionality.
\newblock arXiv:2410.18784, originally posted 2024; revised 2026.

\bibitem[Karatzas and Shreve(1998)]{karatzas1998brownian}
I. Karatzas and S. E. Shreve.
\newblock \emph{Brownian Motion and Stochastic Calculus}.
\newblock Springer, New York, second edition, 1998.

\bibitem[Koike(2026a)]{koike2026follmer}
Y. Koike.
\newblock A note on connections between the F\"ollmer process and the denoising diffusion probabilistic model.
\newblock arXiv:2605.18040, 2026.

\bibitem[Koike(2026b)]{koike2026wasserstein}
Y. Koike.
\newblock Wasserstein bounds for denoising diffusion probabilistic models via the F\"ollmer process.
\newblock arXiv:2605.18069, 2026.

\bibitem[Li and Yan(2024)]{li2024adapting}
G. Li and Y. Yan.
\newblock Adapting to unknown low-dimensional structures in score-based diffusion models.
\newblock In \emph{Advances in Neural Information Processing Systems 37}, 2024.

\bibitem[Li and Yan(2025)]{li2025od}
G. Li and Y. Yan.
\newblock $O(d/T)$ convergence theory for diffusion probabilistic models under minimal assumptions.
\newblock \emph{Journal of Machine Learning Research}, 26:1--55, 2025.

\bibitem[Liang et al.(2025)Liang, Huang, and Chen]{liang2025lowdimtv}
J. Liang, Z. Huang, and Y. Chen.
\newblock Low-dimensional adaptation of diffusion models: convergence in total variation.
\newblock arXiv:2501.12982, accepted at the Conference on Learning Theory, 2025; revised 2026.

\bibitem[Potaptchik et al.(2024)Potaptchik, Azangulov, and Deligiannidis]{potaptchik2024linear}
P. Potaptchik, I. Azangulov, and G. Deligiannidis.
\newblock Linear convergence of diffusion models under the manifold hypothesis.
\newblock arXiv:2410.09046, 2024.

\bibitem[Pope et al.(2021)Pope, Zhu, Abdelkader, Goldblum, and Goldstein]{pope2021intrinsic}
P. Pope, C. Zhu, A. Abdelkader, M. Goldblum, and T. Goldstein.
\newblock The intrinsic dimension of images and its impact on learning.
\newblock In \emph{International Conference on Learning Representations}, 2021.

\bibitem[Song et al.(2021)Song, Sohl-Dickstein, Kingma, Kumar, Ermon, and Poole]{song2021score}
Y. Song, J. Sohl-Dickstein, D. P. Kingma, A. Kumar, S. Ermon, and B. Poole.
\newblock Score-based generative modeling through stochastic differential equations.
\newblock In \emph{International Conference on Learning Representations}, 2021.

\bibitem[Yu and Yu(2025)]{yu2025advancing}
Y. Yu and L. Yu.
\newblock Advancing Wasserstein convergence analysis of score-based models: insights from discretization and second-order acceleration.
\newblock In \emph{Advances in Neural Information Processing Systems 38}, 2025.

\bibitem[Villani(2009)]{villani2009optimal}
C. Villani.
\newblock \emph{Optimal Transport: Old and New}.
\newblock Grundlehren der mathematischen Wissenschaften, volume 338.
Springer, Berlin, Heidelberg, 2009.

\bibitem[Wainwright(2019)]{wainwright2019high}
M. J. Wainwright.
\newblock \emph{High-Dimensional Statistics: A Non-Asymptotic
Viewpoint}.
\newblock Cambridge University Press, Cambridge, 2019.

\end{thebibliography}

\end{document}
